% =====================================================================
\chapter{Soluciones de los ejercicios}
\label{ap:soluciones}
% =====================================================================

Cada solución se presenta en cinco partes: datos, planteamiento, desarrollo, resultado y comprobación. En los ejercicios de demostración, los datos son las hipótesis y el resultado es el enunciado demostrado; la comprobación consiste en un caso particular o en una verificación independiente. Todos los resultados numéricos se han comprobado con un programa, por enumeración exhaustiva cuando el espacio muestral es pequeño.

% ---------------------------------------------------------------------
\section{Tema 1}

\begin{solucion}{eje:1-reglas}
\datos cuatro situaciones de recuento.
\planteamiento «y» entre etapas sucesivas es la regla del producto; «o» entre casos disjuntos, la de la suma; «no todos» sugiere el complemento; los subconjuntos se cuentan por biyección con cadenas de bits.
\pasos
\begin{enumerate}[label=(\alph*)]
  \item Producto: $4\cdot5\cdot3=60$ menús.
  \item Suma (un plato es primero, segundo o postre, casos disjuntos): $4+5+3=12$.
  \item Complemento: $2^{10}-1=1023$ cadenas.
  \item Biyección con $\{0,1\}^{10}$ (\cref{ej:t1-bits}) y producto: $2^{10}=1024$.
\end{enumerate}
\resultado $60$; $12$; $1023$; $1024$.
\comprobacion en (c) y (d), la única cadena excluida en (c) corresponde al subconjunto vacío en (d): $1024-1=1023$.
\end{solucion}

\begin{solucion}{eje:1-vocabulario}
\datos $K=50\,000$, secuencias de longitud 3.
\planteamiento \cref{teo:t1-variaciones}: con repetición $K^{3}$, sin repetición $K^{\underline 3}$; la fracción con algún repetido se obtiene por complemento.
\pasos
\begin{enumerate}[label=\arabic*.]
  \item $K^{3}=50\,000^{3}=1{,}25\cdot10^{14}$.
  \item $K^{\underline3}=50\,000\cdot49\,999\cdot49\,998=124\,992\,500\,100\,000$.
  \item Con algún repetido: $K^{3}-K^{\underline3}=7\,499\,900\,000$.
  \item Fracción: $7\,499\,900\,000/(1{,}25\cdot10^{14})\approx5{,}99992\cdot10^{-5}$.
\end{enumerate}
\resultado (a) $1{,}25\cdot10^{14}$; (b) $124\,992\,500\,100\,000$; (c) $6{,}00\cdot10^{-5}$.
\comprobacion una secuencia tiene algún repetido si coinciden alguno de los 3 pares de posiciones; cada coincidencia tiene proporción $1/K$, y la cota de la suma da $3/K=6\cdot10^{-5}$, muy cerca del valor exacto porque las coincidencias múltiples son rarísimas.
\end{solucion}

\begin{solucion}{eje:1-contrasenas}
\datos longitudes $n\in\{6,7,8\}$, 36 caracteres posibles, de ellos 26 letras.
\planteamiento regla de la suma sobre la longitud (casos disjuntos) y, para cada longitud, complemento: todas menos las que no tienen dígitos.
\pasos
\begin{enumerate}[label=\arabic*.]
  \item $n=6$: $36^{6}-26^{6}=2\,176\,782\,336-308\,915\,776=1\,867\,866\,560$.
  \item $n=7$: $36^{7}-26^{7}=78\,364\,164\,096-8\,031\,810\,176=70\,332\,353\,920$.
  \item $n=8$: $36^{8}-26^{8}=2\,821\,109\,907\,456-208\,827\,064\,576=2\,612\,282\,842\,880$.
  \item Suma: $1\,867\,866\,560+70\,332\,353\,920+2\,612\,282\,842\,880=2\,684\,483\,063\,360$.
\end{enumerate}
\resultado $2\,684\,483\,063\,360\approx2{,}68\cdot10^{12}$ contraseñas.
\comprobacion las de longitud 8 son el $97\,\%$ del total, como cabe esperar, porque cada carácter adicional multiplica el recuento por casi 36.
\end{solucion}

\begin{solucion}{eje:1-divisibles}
\datos $U=\{1,\dots,1000\}$; $A_d$ = múltiplos de $d$, con $\abs{A_d}=\lfloor1000/d\rfloor$.
\planteamiento inclusión--exclusión para tres conjuntos (\cref{teo:t1-ie}); un entero es múltiplo de dos de los números 2, 3, 5 (primos entre sí) si y solo si lo es de su producto.
\pasos
\begin{enumerate}[label=\arabic*.]
  \item $\abs{A_2}=500$, $\abs{A_3}=333$, $\abs{A_5}=200$.
  \item $\abs{A_2\cap A_3}=\abs{A_6}=166$, $\abs{A_2\cap A_5}=\abs{A_{10}}=100$, $\abs{A_3\cap A_5}=\abs{A_{15}}=66$.
  \item $\abs{A_2\cap A_3\cap A_5}=\abs{A_{30}}=33$.
  \item $\abs{A_2\cup A_3\cup A_5}=500+333+200-166-100-66+33=734$.
  \item Ninguno: $1000-734=266$.
\end{enumerate}
\resultado 734 son divisibles por alguno y 266 por ninguno.
\comprobacion la proporción de no divisibles en cada bloque de 30 enteros consecutivos es $\frac12\cdot\frac23\cdot\frac45=\frac{8}{30}$; en 990 enteros hay $33\cdot8=264$, y entre 991 y 1000 solo 991 y 997 no son divisibles por 2, 3 ni 5: $264+2=266$.
\end{solucion}

\begin{solucion}{eje:1-palindromos}
\datos alfabeto de 26 letras.
\planteamiento un palíndromo queda determinado por su primera mitad (redondeando hacia arriba); la aplicación «palíndromo $\mapsto$ primera mitad» es una biyección.
\pasos
\begin{enumerate}[label=\arabic*.]
  \item Longitud 2: la palabra es $xx$; hay 26.
  \item Longitud 3: la palabra es $xyx$; hay $26^{2}=676$.
  \item Proporciones: $26/26^{2}=1/26$ y $26^{2}/26^{3}=1/26$: son iguales.
  \item Longitud $n$: se eligen libremente las $\lceil n/2\rceil$ primeras letras; hay $26^{\lceil n/2\rceil}$.
\end{enumerate}
\resultado 26 y 676; ambas proporciones valen $1/26$; en general, $26^{\lceil n/2\rceil}$.
\comprobacion la proporción $1/26$ en los dos casos se explica directamente: una palabra de longitud 2 o 3 es palíndromo si y solo si su última letra coincide con la primera.
\end{solucion}

\begin{solucion}{eje:1-mesa}
\datos 8 personas, disposiciones circulares módulo rotación.
\planteamiento \cref{ej:t1-circular}; para (b), se pega a Ana y Blas en un bloque.
\pasos
\begin{enumerate}[label=\arabic*.]
  \item (a) $(8-1)!=7!=5040$.
  \item (b) El bloque Ana--Blas y las otras 6 personas forman 7 unidades: $(7-1)!=720$ disposiciones circulares. Dentro del bloque hay 2 órdenes (Ana a la izquierda o a la derecha de Blas): $2\cdot720=1440$.
  \item (c) Complemento: $5040-1440=3600$.
\end{enumerate}
\resultado 5040; 1440; 3600.
\comprobacion fijada Ana, Blas ocupa uno de los 7 asientos restantes con la misma frecuencia, y 2 de ellos son contiguos a Ana: $5040\cdot\frac27=1440$.
\end{solucion}

\begin{solucion}{eje:1-cifras}
\datos números $\overline{abcd}$ con $a\neq0$ y cifras distintas.
\planteamiento regla del producto generalizada, eligiendo primero las cifras con restricciones; para los pares se separan casos según la última cifra sea 0 o no.
\pasos
\begin{enumerate}[label=\arabic*.]
  \item Total: $a$ admite 9 valores; $b$, 9 (cualquiera salvo $a$, incluido el 0); $c$, 8; $d$, 7: $9\cdot9\cdot8\cdot7=4536$.
  \item Pares con $d=0$: $a$ admite 9 valores, $b$ 8 y $c$ 7: $504$.
  \item Pares con $d\in\{2,4,6,8\}$: $d$ admite 4 valores; $a$, 8 (ni 0 ni $d$); $b$, 8; $c$, 7: $4\cdot8\cdot8\cdot7=1792$.
  \item Pares: $504+1792=2296$.
\end{enumerate}
\resultado 4536 números con cifras distintas, de los que 2296 son pares.
\comprobacion enumeración exhaustiva con un programa: 4536 y 2296.
\end{solucion}

\begin{solucion}{eje:1-funciones}
\datos $\abs{A}=4$, $\abs{B}=2$; en (d), conjuntos de 5 y 3 elementos.
\planteamiento \cref{teo:t1-variaciones}; una función no es sobreyectiva si su imagen omite algún elemento del codominio.
\pasos
\begin{enumerate}[label=(\alph*)]
  \item $2^{4}=16$.
  \item $2^{\underline4}=2\cdot1\cdot0\cdot(-1)=0$: no hay ninguna (palomar).
  \item No sobreyectivas: las constantes, 2. Sobreyectivas: $16-2=14$.
  \item Sea $E_y$ el conjunto de funciones que no alcanzan $y\in B$, $\abs{B}=3$. $\abs{E_y}=2^{5}=32$, $\abs{E_y\cap E_{y'}}=1$ y la intersección de los tres es vacía. No sobreyectivas: $3\cdot32-3\cdot1+0=93$. Sobreyectivas: $3^{5}-93=243-93=150$.
\end{enumerate}
\resultado 16; 0; 14; 150.
\comprobacion enumeración exhaustiva con un programa: 14 sobreyecciones $4\to2$ y 150 sobreyecciones $5\to3$. Además, $150=3!\cdot25$, donde 25 es el número de formas de partir un conjunto de 5 elementos en 3 bloques no vacíos: cada partición, con sus bloques asignados a los 3 elementos de $B$, da una sobreyección.
\end{solucion}

\begin{solucion}{eje:1-sombreros}
\datos $n$ personas, repartos = permutaciones de los $n$ sombreros.
\planteamiento regla del producto sobre las personas, contando cuántos sombreros admite cada una.
\pasos
\begin{enumerate}[label=(\alph*)]
  \item Solo la permutación identidad: 1 reparto.
  \item Las $m$ primeras tienen su sombrero fijado; los $n-m$ restantes se reparten libremente: $(n-m)!$.
  \item Las $m$ primeras personas reciben, en algún orden, los $m$ sombreros de las $m$ últimas: $m!$ formas (biyecciones entre dos conjuntos de $m$ elementos). Los otros $n-m$ sombreros se reparten entre las $n-m$ personas restantes: $(n-m)!$. Total, $m!\,(n-m)!$.
\end{enumerate}
\resultado $1$; $(n-m)!$; $m!\,(n-m)!$.
\comprobacion con $n=4$, $m=2$: (b) da $2$ y (c) da $2!\cdot2!=4$; enumerando las 24 permutaciones de $\{1,2,3,4\}$ se obtienen 2 con $\sigma(1)=1$, $\sigma(2)=2$ y 4 con $\{\sigma(1),\sigma(2)\}=\{3,4\}$.
\end{solucion}

\begin{solucion}{eje:1-encuesta}
\datos $\abs{U}=60$, $\abs{C}=30$, $\abs{A}=25$, $\abs{P}=20$, $\abs{C\cap A}=10$, $\abs{C\cap P}=8$, $\abs{A\cap P}=7$, $\abs{C\cap A\cap P}=3$.
\planteamiento inclusión--exclusión para la unión; para «exactamente una», se resta a cada asignatura lo compartido.
\pasos
\begin{enumerate}[label=\arabic*.]
  \item $\abs{C\cup A\cup P}=30+25+20-10-8-7+3=53$; ninguna: $60-53=7$.
  \item Solo Cálculo: $30-10-8+3=15$; solo Álgebra: $25-10-7+3=11$; solo Programación: $20-8-7+3=8$.
  \item Exactamente una: $15+11+8=34$.
\end{enumerate}
\resultado 7 no cursan ninguna y 34 cursan exactamente una.
\comprobacion exactamente dos: $(10-3)+(8-3)+(7-3)=16$; exactamente tres: 3. Total en la unión: $34+16+3=53$.
\end{solucion}

\begin{solucion}{eje:1-producto-ind}
\datos $\abs{B}=n$; $I(k)$ = número de inyecciones $[k]\to B$.
\planteamiento inducción sobre $k$, descomponiendo cada inyección en su restricción a $[k]$ y el valor en $k+1$, y usando la regla de la suma.
\pasos
\begin{enumerate}[label=\arabic*.]
  \item Caso base, $k=0$: hay una única función $\emptyset\to B$ (la vacía), que es inyectiva; $n^{\underline0}=1$.
  \item Paso: sea $0\le k<n$ y supóngase $I(k)=n^{\underline k}$. Si $f\colon[k+1]\to B$ es inyectiva, $g=f|_{[k]}$ es inyectiva y $f(k+1)\notin g([k])$. Recíprocamente, cada par $(g,b)$ con $g$ inyectiva y $b\in B\setminus g([k])$ define una única inyección $f$. Es una biyección.
  \item Para cada $g$, $\abs{B\setminus g([k])}=n-k$, porque $g$ es inyectiva y $\abs{g([k])}=k$. Por la regla de la suma sobre las $g$, $I(k+1)=I(k)(n-k)=n^{\underline k}(n-k)=n^{\underline{k+1}}$.
  \item Si $k>n$: la restricción a $[n+1]$ de una inyección $[k]\to B$ sería una inyección $[n+1]\to B$, y el paso anterior con $k=n$ da $I(n+1)=n^{\underline n}\cdot(n-n)=0$.
\end{enumerate}
\resultado hay $n^{\underline k}$ inyecciones si $k\le n$ y ninguna si $k>n$: no se pueden colocar $k$ palomas en $n<k$ nidos sin que dos compartan nido.
\comprobacion $n=3$, $k=2$: las inyecciones $[2]\to\{a,b,c\}$ son los 6 pares ordenados de elementos distintos, y $3^{\underline2}=6$.
\end{solucion}

\begin{solucion}{eje:1-fibonacci}
\datos $a_n$ = número de cadenas de $n$ bits sin \texttt{11}.
\planteamiento clasificar las cadenas válidas según su último bit.
\pasos
\begin{enumerate}[label=\arabic*.]
  \item $a_1=2$ (\texttt{0}, \texttt{1}); $a_2=3$ (\texttt{00}, \texttt{01}, \texttt{10}).
  \item Sea $n\ge3$ y $S_0$, $S_1$ las cadenas válidas de longitud $n$ que terminan en \texttt{0} y en \texttt{1}. Son disjuntas y su unión es el conjunto de cadenas válidas (regla de la suma).
  \item $w\mapsto w\texttt{0}$ es una biyección entre las cadenas válidas de longitud $n-1$ y $S_0$: añadir un \texttt{0} no crea \texttt{11}, y quitar el \texttt{0} final de una cadena válida da una válida. Luego $\abs{S_0}=a_{n-1}$.
  \item Si una cadena válida termina en \texttt{1}, su penúltimo bit es \texttt{0}. Por tanto $w\mapsto w\texttt{01}$ es una biyección entre las cadenas válidas de longitud $n-2$ y $S_1$. Luego $\abs{S_1}=a_{n-2}$.
  \item $a_n=a_{n-1}+a_{n-2}$; sucesivamente $a_3=5$, $a_4=8$, $a_5=13$, $a_6=21$, $a_7=34$, $a_8=55$.
\end{enumerate}
\resultado la recurrencia de Fibonacci con $a_1=2$, $a_2=3$, y $a_8=55$.
\comprobacion enumeración exhaustiva de las 256 cadenas de 8 bits: 55 no contienen \texttt{11}.
\end{solucion}

\begin{solucion}{eje:1-ie-cuatro}
\datos $A,B,C,D$ finitos; \cref{teo:t1-ie}.
\planteamiento aplicar la fórmula de dos conjuntos a $X=A\cup B\cup C$ y $D$, y la de tres conjuntos dos veces.
\pasos
\begin{enumerate}[label=\arabic*.]
  \item $\abs{X\cup D}=\abs X+\abs D-\abs{X\cap D}$.
  \item Por la distributiva, $X\cap D=(A\cap D)\cup(B\cap D)\cup(C\cap D)$, y por la fórmula de tres conjuntos, con $(A\cap D)\cap(B\cap D)=A\cap B\cap D$, etc.:
  \[
  \abs{X\cap D}=\abs{A\cap D}+\abs{B\cap D}+\abs{C\cap D}-\abs{A\cap B\cap D}-\abs{A\cap C\cap D}-\abs{B\cap C\cap D}+\abs{A\cap B\cap C\cap D}.
  \]
  \item $\abs X$ se desarrolla con la fórmula de tres conjuntos.
  \item Sustituyendo: suma de los 4 cardinales, menos los $\binom42=6$ de las intersecciones dobles ($AB$, $AC$, $BC$ de $\abs X$ y $AD$, $BD$, $CD$ de $\abs{X\cap D}$), más los $\binom43=4$ de las triples ($ABC$ de $\abs X$ y $ABD$, $ACD$, $BCD$), menos la cuádruple.
\end{enumerate}
\resultado $\abs{A\cup B\cup C\cup D}=\sum\abs{A_i}-\sum_{i<j}\abs{A_i\cap A_j}+\sum_{i<j<k}\abs{A_i\cap A_j\cap A_k}-\abs{A\cap B\cap C\cap D}$.
\comprobacion con $A=B=C=D=\{1\}$: $4-6+4-1=1=\abs{\{1\}}$.
\end{solucion}

\begin{solucion}{eje:1-collares}
\datos (a) $n\ge3$ personas en una mesa redonda, módulo rotación y reflexión; (b) $2n$ personas en $n$ parejas.
\planteamiento regla de la división en ambos casos, identificando cuántos objetos «ordenados» dan el mismo objeto «sin orden».
\pasos
\begin{enumerate}[label=\arabic*.]
  \item (a) Sea $f$ la aplicación que lleva cada disposición circular (de las $(n-1)!$) a su clase módulo reflexión. Cada clase contiene una disposición y su reflejada.
  \item Ambas son distintas si $n\ge3$: en la reflejada, el vecino de la derecha de una persona $x$ es el que en la original era su vecino de la izquierda, y para $n\ge3$ los dos vecinos de $x$ son personas distintas.
  \item $f$ es 2 a~1, luego hay $(n-1)!/2$ clases.
  \item (b) Sea $g$ la aplicación que a cada permutación $(a_1,\dots,a_{2n})$ asocia el emparejamiento $\{\{a_1,a_2\},\dots,\{a_{2n-1},a_{2n}\}\}$; es sobreyectiva.
  \item Un emparejamiento dado se obtiene de las permutaciones que ordenan sus $n$ parejas de cualquier forma ($n!$) y, dentro de cada pareja, ponen primero cualquiera de los dos ($2^{n}$): $g$ es $2^{n}n!$ a~1.
  \item Por la regla de la división, hay $(2n)!/(2^{n}n!)$ emparejamientos.
\end{enumerate}
\resultado $(n-1)!/2$ y $(2n)!/(2^{n}n!)$.
\comprobacion (a) con $n=3$: $2!/2=1$, pues tres personas en círculo, sin distinguir sentido, están siempre en la misma disposición. (b) con $n=2$: $4!/(4\cdot2)=3$, y en efecto $\{a,b,c,d\}$ admite los emparejamientos $ab|cd$, $ac|bd$, $ad|bc$.
\end{solucion}

% ---------------------------------------------------------------------
\section{Tema 2}

\begin{solucion}{eje:2-calcular}
\planteamiento \cref{teo:t2-binom} y \cref{teo:t2-binomio}.
\pasos
\begin{enumerate}[label=\arabic*.]
  \item $\binom{10}{3}=\frac{10\cdot9\cdot8}{3!}=120$ y $\binom{10}{7}=\frac{10!}{7!\,3!}=120$.
  \item $\binom{52}{5}=\frac{52\cdot51\cdot50\cdot49\cdot48}{120}=2\,598\,960$.
  \item El coeficiente de $x^{3}$ en $(1+x)^{7}$ es $\binom73=35$.
  \item Elegir los 3 elementos que \emph{entran} equivale a elegir los 7 que \emph{quedan fuera} (\cref{prop:t2-props}(i)).
\end{enumerate}
\resultado $120$, $120$, $2\,598\,960$ y $35$.
\comprobacion fila 7 del triángulo de Pascal: $1,7,21,35,35,21,7,1$, que suma $128=2^{7}$.
\end{solucion}

\begin{solucion}{eje:2-modelo}
\planteamiento decidir si importa el orden y si se permite repetir (\cref{tab:t2-modelos}).
\pasos
\begin{enumerate}[label=(\alph*)]
  \item Ordenado, sin repetición: $8^{\underline3}=336$.
  \item No ordenado, sin repetición: $\binom83=56$.
  \item Ordenado, con repetición: $8^{3}=512$.
  \item No ordenado, con repetición: $\binom{8+3-1}{3}=\binom{10}{3}=120$.
\end{enumerate}
\resultado $336$, $56$, $512$ y $120$.
\comprobacion $336=56\cdot3!$: cada comisión da $3!$ podios. En cambio $512/3!$ no es entero, lo que confirma que (d) no puede obtenerse dividiendo (c) por $3!$.
\end{solucion}

\begin{solucion}{eje:2-poker}
\datos 52 cartas, 13 valores y 4 palos; $\binom{52}5=2\,598\,960$ manos.
\planteamiento construir cada mano por etapas (valores, después palos) de modo que cada mano se obtenga exactamente una vez~\cite[§14.7]{mcs}.
\pasos
\begin{enumerate}[label=(\alph*)]
  \item Póquer: valor del cuarteto (13) y quinta carta (48 restantes): $13\cdot48=624$.
  \item Full: valor del trío (13), sus palos $\binom43=4$, valor de la pareja (12), sus palos $\binom42=6$: $13\cdot4\cdot12\cdot6=3744$.
  \item Doble pareja: los dos valores de las parejas, sin orden, $\binom{13}{2}=78$; palos de cada pareja, $6\cdot6=36$; quinta carta de otro valor, $52-8=44$: $78\cdot36\cdot44=123\,552$.
  \item Fracciones: $624/2\,598\,960=1/4165\approx0{,}00024$; $3744/2\,598\,960=6/4165\approx0{,}00144$; $123\,552/2\,598\,960=198/4165\approx0{,}0475$.
\end{enumerate}
\resultado 624, 3744 y 123\,552 manos.
\comprobacion si en (c) se eligieran los valores de las parejas en orden ($13\cdot12$) se obtendría el doble, $247\,104$: cada mano se contaría dos veces.
\end{solucion}

\begin{solucion}{eje:2-equipos}
\datos 12 personas.
\planteamiento \cref{teo:t2-multinomial} para equipos con nombre y regla de la división para los equipos del mismo tamaño, que son indistinguibles.
\pasos
\begin{enumerate}[label=\arabic*.]
  \item (a) Con nombres (el de 2, el «primero» de 5 y el «segundo» de 5): $\binom{12}{2,5,5}=\frac{12!}{2!\,5!\,5!}=16\,632$.
  \item Los dos equipos de 5 no tienen nombre: cada reparto se ha contado $2!$ veces. Resultado: $16\,632/2=8316$.
  \item (b) Con nombres: $\frac{12!}{4!\,4!\,4!}=34\,650$; sin nombres, se divide por $3!$: $5775$.
\end{enumerate}
\resultado (a) 8316; (b) 5775.
\comprobacion (b) por etapas: la persona~1 elige 3 compañeros entre 11, $\binom{11}{3}=165$; la persona de menor número entre las restantes elige 3 entre 7, $\binom73=35$; el resto forma el tercer equipo: $165\cdot35=5775$.
\end{solucion}

\begin{solucion}{eje:2-anagramas}
\datos \texttt{MISSISSIPPI}: M~1, I~4, S~4, P~2 (11 letras). \texttt{PROBABILIDAD}: B, A, I, D dos veces cada una y P, R, O, L una vez (12 letras).
\planteamiento \cref{teo:t2-multinomial}; para las \texttt{S} juntas, se trata el bloque \texttt{SSSS} como un solo símbolo.
\pasos
\begin{enumerate}[label=\arabic*.]
  \item $\frac{11!}{1!\,4!\,4!\,2!}=\frac{39\,916\,800}{1152}=34\,650$.
  \item $\frac{12!}{2!^{4}}=\frac{479\,001\,600}{16}=29\,937\,600$.
  \item Con el bloque: símbolos M, I~$\times4$, P~$\times2$ y el bloque: 8 símbolos, $\frac{8!}{4!\,2!}=840$.
\end{enumerate}
\resultado 34\,650; 29\,937\,600; 840.
\comprobacion la proporción de anagramas con las \texttt{S} juntas, $840/34\,650\approx0{,}024$, coincide con la probabilidad de que las 4 posiciones de las \texttt{S}, elegidas al azar entre 11, sean consecutivas: $8/\binom{11}{4}=8/330\approx0{,}024$.
\end{solucion}

\begin{solucion}{eje:2-coeficientes}
\planteamiento binomio con $x\to2x$, $y\to-y$; multinomio de Newton; los términos de $(x+y+z)^{6}$ corresponden a las soluciones de $k_1+k_2+k_3=6$.
\pasos
\begin{enumerate}[label=\arabic*.]
  \item Término general de $(2x-y)^{8}$: $\binom8k(2x)^{k}(-y)^{8-k}$. Con $k=3$: $\binom83\,2^{3}(-1)^{5}=56\cdot8\cdot(-1)=-448$.
  \item $\binom{6}{2,3,1}=\frac{720}{2\cdot6\cdot1}=60$.
  \item Número de términos: soluciones de $k_1+k_2+k_3=6$ con $k_i\ge0$, $\binom{8}{2}=28$ (\cref{teo:t2-estrellas}).
\end{enumerate}
\resultado $-448$; $60$; 28 términos.
\comprobacion la suma de los coeficientes de $(x+y+z)^{6}$ es $3^{6}=729$ (tomar $x=y=z=1$), compatible con 28 términos cuyo mayor coeficiente es $\binom{6}{2,2,2}=90$.
\end{solucion}

\begin{solucion}{eje:2-ecuaciones}
\planteamiento \cref{teo:t2-estrellas} con $n=4$ incógnitas y cambio de variable para las cotas inferiores.
\pasos
\begin{enumerate}[label=(\alph*)]
  \item $\binom{10+3}{3}=\binom{13}3=286$.
  \item $\binom{10-1}{3}=\binom93=84$.
  \item Con $y_1=x_1-2\ge0$: $y_1+x_2+x_3+x_4=8$, $\binom{11}{3}=165$.
\end{enumerate}
\resultado 286, 84 y 165.
\comprobacion por complemento, las soluciones de (a) con $x_1\le1$ son las de $x_2+x_3+x_4=10$ y las de $x_2+x_3+x_4=9$: $\binom{12}2+\binom{11}2=66+55=121$, y $286-121=165$.
\end{solucion}

\begin{solucion}{eje:2-caminos}
\planteamiento \cref{ej:t2-caminos}: de $(0,0)$ a $(a,b)$ hay $\binom{a+b}{a}$ caminos; pasar por un punto intermedio divide el camino en dos tramos independientes (regla del producto).
\pasos
\begin{enumerate}[label=\arabic*.]
  \item (a) $\binom{11}{5}=462$.
  \item (b) De $(0,0)$ a $(2,3)$: $\binom52=10$; de $(2,3)$ a $(5,6)$, es decir, 3 a la derecha y 3 arriba: $\binom63=20$. Total, $200$.
  \item (c) $462-200=262$.
\end{enumerate}
\resultado 462; 200; 262.
\comprobacion la proporción $200/462\approx0{,}43$ es la probabilidad de que un camino uniforme pase por $(2,3)$; es razonable que sea alta porque $(2,3)$ está cerca de la diagonal que va a $(5,6)$.
\end{solucion}

\begin{solucion}{eje:2-estaciones}
\datos 7 personas, 4 estaciones; $4^{7}=16\,384$ asignaciones.
\planteamiento complemento e inclusión--exclusión (\cref{teo:t2-ie}) sobre los sucesos $E_s$ = «nadie nace en la estación $s$».
\pasos
\begin{enumerate}[label=\arabic*.]
  \item $\abs{E_s}=3^{7}=2187$; $\abs{E_s\cap E_t}=2^{7}=128$; tres estaciones vacías: $1^{7}=1$; cuatro: 0.
  \item $\abs{\bigcup E_s}=4\cdot2187-6\cdot128+4\cdot1-0=8748-768+4=7984$.
  \item Con las cuatro estaciones: $16\,384-7984=8400$.
  \item Fracción: $8400/16\,384\approx0{,}5127$.
\end{enumerate}
\resultado 8400 asignaciones, un $51{,}3\,\%$ del total.
\comprobacion es el número de sobreyecciones de un conjunto de 7 en uno de 4, que una enumeración con un programa confirma: 8400.
\end{solucion}

\begin{solucion}{eje:2-desarreglos}
\datos $n=5$.
\planteamiento \cref{ej:t2-desarreglos}.
\pasos
\begin{enumerate}[label=\arabic*.]
  \item $D_5=5!\bigl(1-1+\tfrac12-\tfrac16+\tfrac1{24}-\tfrac1{120}\bigr)=60-20+5-1=44$.
  \item Fracción: $44/120\approx0{,}3667$; $1/e\approx0{,}3679$.
\end{enumerate}
\resultado 44 sorteos válidos, un $36{,}7\,\%$, a menos de $0{,}002$ de $1/e$.
\comprobacion recurrencia $D_n=(n-1)(D_{n-1}+D_{n-2})$ con $D_1=0$, $D_2=1$: $D_3=2$, $D_4=9$, $D_5=4(9+2)=44$.
\end{solucion}

\begin{solucion}{eje:2-ositos}
\datos 5 sabores, $n$ ositos con $30\le n\le50$.
\planteamiento para cada $n$, una bolsa es una solución de $x_1+\dots+x_5=n$ con $x_i\ge0$: $\binom{n+4}{4}$ bolsas. Se suma sobre $n$ y se usa la identidad del palo de \emph{hockey} (\cref{eje:2-palo}).
\pasos
\begin{enumerate}[label=\arabic*.]
  \item Total: $\sum_{n=30}^{50}\binom{n+4}{4}=\sum_{j=34}^{54}\binom j4$.
  \item Por el \cref{eje:2-palo}, $\sum_{j=4}^{54}\binom j4=\binom{55}5$ y $\sum_{j=4}^{33}\binom j4=\binom{34}5$.
  \item Restando: $\binom{55}{5}-\binom{34}{5}=3\,478\,761-278\,256=3\,200\,505$.
\end{enumerate}
\resultado $\binom{55}5-\binom{34}5=3\,200\,505$ bolsas.
\comprobacion otra lectura: añadiendo un sexto «sabor» ficticio $x_6=50-n$, las bolsas con a lo sumo 50 ositos son las soluciones de $x_1+\dots+x_6=50$, $\binom{55}{5}$, y las de a lo sumo 29 son $\binom{34}{5}$; la suma directa con un programa da el mismo valor.
\end{solucion}

\begin{solucion}{eje:2-palo}
\datos enteros $n\ge k\ge0$.
\planteamiento historia: contar los grupos de $k+1$ personas elegidas entre $n+1$ personas de edades distintas, clasificándolos por la persona de mayor edad del grupo.
\pasos
\begin{enumerate}[label=\arabic*.]
  \item Se numeran las personas $1,\dots,n+1$ de menor a mayor edad. El lado derecho, $\binom{n+1}{k+1}$, es el número de grupos de $k+1$.
  \item Si la persona de más edad del grupo es la $j+1$, los otros $k$ miembros se eligen entre las $j$ personas más jóvenes que ella: $\binom jk$ grupos. Para que sea posible, $j\ge k$, y como $j+1\le n+1$, $j\le n$.
  \item Cada grupo tiene exactamente una persona de mayor edad, así que las clases $j=k,\dots,n$ son disjuntas y cubren todos los grupos. Por la regla de la suma, $\binom{n+1}{k+1}=\sum_{j=k}^{n}\binom jk$.
\end{enumerate}
\resultado $\sum_{j=k}^{n}\binom jk=\binom{n+1}{k+1}$.
\comprobacion $k=1$, $n=4$: $1+2+3+4=10=\binom52$.
\end{solucion}

\begin{solucion}{eje:2-sumak}
\datos $n\ge1$ (en (b), $n\ge2$).
\planteamiento (a) absorción y \cref{cor:t2-sumas}; (b) doble conteo de comités con presidente y vicepresidente.
\pasos
\begin{enumerate}[label=\arabic*.]
  \item (a) Para $k\ge1$, $k\binom nk=n\binom{n-1}{k-1}$; el término $k=0$ es nulo. Luego $\sum_kk\binom nk=n\sum_{k=1}^{n}\binom{n-1}{k-1}=n\,2^{n-1}$.
  \item (b) Se cuentan las ternas (comité $C\subseteq[n]$, presidente $p\in C$, vicepresidente $v\in C$, $v\neq p$). Por tamaño del comité: si $\abs C=k$, hay $\binom nk$ comités y $k(k-1)$ elecciones de $(p,v)$; total $\sum_{k=2}^{n}k(k-1)\binom nk$.
  \item Eligiendo primero los cargos: $n$ presidentes, $n-1$ vicepresidentes, y el resto del comité es un subconjunto cualquiera de las $n-2$ personas restantes, $2^{n-2}$: total $n(n-1)2^{n-2}$.
\end{enumerate}
\resultado $\sum_kk\binom nk=n2^{n-1}$ y $\sum_kk(k-1)\binom nk=n(n-1)2^{n-2}$.
\comprobacion $n=3$: $1\cdot3+2\cdot3+3\cdot1=12=3\cdot2^{2}$ y $2\cdot1\cdot3+3\cdot2\cdot1=12=3\cdot2\cdot2^{1}$.
\end{solucion}

\begin{solucion}{eje:2-vandermonde}
\datos $m$ mujeres, $n$ hombres, comités de $k$ personas.
\planteamiento doble conteo: total y clasificación por el número $j$ de mujeres.
\pasos
\begin{enumerate}[label=\arabic*.]
  \item Hay $\binom{m+n}{k}$ comités de $k$ personas.
  \item Los comités con exactamente $j$ mujeres se forman eligiendo $j$ mujeres y $k-j$ hombres: $\binom mj\binom n{k-j}$ (con el convenio $\binom ab=0$ si $b>a$ o $b<0$).
  \item Las clases $j=0,\dots,k$ son disjuntas y lo cubren todo; por la regla de la suma, $\binom{m+n}k=\sum_j\binom mj\binom n{k-j}$.
  \item Con $m=n=k$: $\binom{2n}{n}=\sum_j\binom nj\binom n{n-j}=\sum_j\binom nj^{2}$, por la simetría.
\end{enumerate}
\resultado la identidad de Vandermonde y $\sum_j\binom nj^{2}=\binom{2n}n$.
\comprobacion $n=3$: $1+9+9+1=20=\binom63$.
\end{solucion}

\begin{solucion}{eje:2-pares}
\datos $\abs A=n\ge1$, $a\in A$; $\varphi(X)=X\dsim\{a\}$.
\planteamiento probar que $\varphi$ es una involución que cambia la paridad del cardinal.
\pasos
\begin{enumerate}[label=\arabic*.]
  \item Si $a\notin X$, $\varphi(X)=X\cup\{a\}$ tiene $\abs X+1$ elementos; si $a\in X$, $\varphi(X)=X\setminus\{a\}$ tiene $\abs X-1$. En ambos casos cambia la paridad.
  \item $\varphi(\varphi(X))=X$: añadir y quitar $a$ (o al revés) devuelve $X$. Luego $\varphi$ es biyectiva, con inversa ella misma.
  \item Restringida a los subconjuntos de cardinal par, $\varphi$ toma valores en los de cardinal impar y es biyectiva sobre ellos (su inversa es $\varphi$ restringida a los impares). Hay tantos subconjuntos pares como impares.
  \item $\sum_k(-1)^{k}\binom nk=\#\text{pares}-\#\text{impares}=0$. Como $\#\text{pares}+\#\text{impares}=2^{n}$, hay $2^{n-1}$ de cada clase.
\end{enumerate}
\resultado $\sum_k(-1)^{k}\binom nk=0$ y hay $2^{n-1}$ subconjuntos de cardinal par.
\comprobacion $n=3$: pares $\emptyset$, $\{a,b\}$, $\{a,c\}$, $\{b,c\}$, que son 4 $=2^{2}$.
\end{solucion}

% ---------------------------------------------------------------------
\section{Tema 3}

\begin{solucion}{eje:3-espacio}
\planteamiento elegir resultados lo bastante detallados para describir el experimento y, si es posible, equiprobables.
\pasos
\begin{enumerate}[label=(\alph*)]
  \item $\Omega=\{\mathrm C,\mathrm X\}^{3}$, con $2^{3}=8$ resultados; Laplace es razonable si la moneda es equilibrada.
  \item $\Omega$ = subconjuntos de 5 cartas, $\binom{52}5=2\,598\,960$; Laplace es razonable si la baraja está bien mezclada.
  \item $\Omega=\V^{2}$, con $K^{2}$ resultados; Laplace \emph{no} es razonable: el modelo asigna probabilidades muy distintas a unas secuencias y otras.
  \item $\Omega=\{1,2,3,\dots\}$, infinito numerable; Laplace no es aplicable (no hay ley uniforme sobre un conjunto infinito numerable, \cref{eje:3-laplace}), y además $\PP(\{k\})=2^{-k}$ con una moneda equilibrada.
\end{enumerate}
\resultado $8$; $2\,598\,960$; $K^{2}$; infinito numerable. Laplace solo en (a) y (b).
\comprobacion en (d), $\sum_{k\ge1}2^{-k}=1$, de modo que la ley indicada es válida (\cref{prop:t3-discreto}).
\end{solucion}

\begin{solucion}{eje:3-sucesos}
\planteamiento traducir «o», «y» y «no» en unión, intersección y complementario; «como máximo uno» es el complementario de «al menos dos».
\pasos
\begin{enumerate}[label=(\alph*)]
  \item $A\cup B\cup C$.
  \item $A\cap B\cap C$.
  \item $A^{c}\cap B^{c}\cap C^{c}=(A\cup B\cup C)^{c}$.
  \item $(A\cap B^{c}\cap C^{c})\cup(A^{c}\cap B\cap C^{c})\cup(A^{c}\cap B^{c}\cap C)$.
  \item $(A^{c}\cap B^{c})\cup(A^{c}\cap C^{c})\cup(B^{c}\cap C^{c})$: entre cada dos, al menos uno no ocurre.
  \item $A\cap B\cap C^{c}$.
\end{enumerate}
\resultado las expresiones anteriores.
\comprobacion (e) es el complementario de «al menos dos», $(A\cap B)\cup(A\cap C)\cup(B\cap C)$, por las leyes de De Morgan; y es la unión disjunta de (c) y (d).
\end{solucion}

\begin{solucion}{eje:3-monedas}
\datos $\Omega=\{\mathrm C,\mathrm X\}^{3}$, 8 resultados equiprobables.
\planteamiento modelo de Laplace.
\pasos
\begin{enumerate}[label=(\alph*)]
  \item $\{\mathrm{CCC}\}$: $1/8$.
  \item $\{\mathrm{CXC}\}$: $1/8$.
  \item Exactamente dos caras: $\binom32=3$ resultados, $3/8$.
  \item Al menos tantas caras como cruces equivale a 2 o 3 caras: $3+1=4$ resultados, $1/2$.
\end{enumerate}
\resultado $1/8$, $1/8$, $3/8$, $1/2$.
\comprobacion (d) vale $1/2$ por simetría: intercambiar caras y cruces lleva «más caras» a «más cruces», y con 3 lanzamientos no puede haber empate.
\end{solucion}

\begin{solucion}{eje:3-dados}
\datos 36 resultados equiprobables $(i,j)$.
\planteamiento Laplace; complemento en (b); en (c), $\{\max=4\}=\{\max\le4\}\setminus\{\max\le3\}$.
\pasos
\begin{enumerate}[label=(\alph*)]
  \item Suma 7: 6 pares, $6/36=1/6$.
  \item Ningún 6: $5^{2}=25$; al menos un 6: $1-25/36=11/36$.
  \item $\abs{\{\max\le4\}}=16$, $\abs{\{\max\le3\}}=9$: $(16-9)/36=7/36$.
\end{enumerate}
\resultado $1/6$, $11/36$, $7/36$.
\comprobacion $\frac16+\frac16=\frac{12}{36}$ cuenta dos veces el resultado $(6,6)$, que está en ambos sucesos: $\frac{12}{36}-\frac1{36}=\frac{11}{36}$ (\cref{prop:t3-props}(vi)).
\end{solucion}

\begin{solucion}{eje:3-cuatro-dados}
\datos 4 dados, $6^{4}=1296$ resultados.
\planteamiento la aplicación $(x_1,\dots,x_4)\mapsto(7-x_1,\dots,7-x_4)$ es una biyección de $\Omega$ que lleva suma $s$ a suma $28-s$.
\pasos
\begin{enumerate}[label=\arabic*.]
  \item Por la biyección, $\PP(S=21)=\PP(S=7)$ y $\PP(S=22)=\PP(S=6)$.
  \item Sumas pequeñas: con $y_i=x_i-1\ge0$, la suma 7 equivale a $y_1+\dots+y_4=3$, con $\binom63=20$ soluciones (todas con $y_i\le5$); la suma 6 equivale a $y_1+\dots+y_4=2$, con $\binom53=10$.
  \item $\PP(S=21)=20/1296\approx0{,}0154$ y $\PP(S=22)=10/1296\approx0{,}0077$.
\end{enumerate}
\resultado es mayor la probabilidad de 21, el doble que la de 22.
\comprobacion enumeración de los 1296 resultados: 20 con suma 21 y 10 con suma 22.
\end{solucion}

\begin{solucion}{eje:3-horarios}
\datos 3 asignaturas, 10 franjas; $10^{3}$ asignaciones equiprobables.
\planteamiento complemento: todas en franjas distintas.
\pasos
\begin{enumerate}[label=\arabic*.]
  \item Franjas distintas: $10\cdot9\cdot8=720$.
  \item $\PP(\text{coincidencia})=1-720/1000=0{,}28$.
\end{enumerate}
\resultado $7/25=0{,}28$.
\comprobacion por la cota de la unión sobre los 3 pares, $\PP\le3/10$; el valor exacto es algo menor porque se descuentan las coincidencias triples.
\end{solucion}

\begin{solucion}{eje:3-robos}
\datos 6 robos, 6 distritos, $6^{6}=46\,656$ asignaciones equiprobables.
\planteamiento complemento: ningún distrito repite, es decir, la asignación es una biyección.
\pasos
\begin{enumerate}[label=\arabic*.]
  \item Asignaciones sin repetición: $6!=720$.
  \item $\PP=1-720/46\,656=1-5/324=319/324$.
\end{enumerate}
\resultado $319/324\approx0{,}9846$.
\comprobacion es el problema del cumpleaños (\cref{prop:t3-cumple}) con $N=n=6$: $1-\prod_{i=0}^{5}(1-i/6)=1-\frac{6!}{6^{6}}$.
\end{solucion}

\begin{solucion}{eje:3-seis-cartas}
\datos manos de 6 cartas, $\binom{52}6=20\,358\,520$.
\planteamiento construir cada mano por valores y palos, como en el \cref{eje:2-poker}.
\pasos
\begin{enumerate}[label=\arabic*.]
  \item Doble pareja: valores de las parejas $\binom{13}2=78$; palos $\binom42^{2}=36$; valores de las dos cartas sueltas, distintos entre sí y de los anteriores, $\binom{11}2=55$; palos de esas dos, $4^{2}=16$. Total: $78\cdot36\cdot55\cdot16=2\,471\,040$.
  \item Trío: valor del trío 13, palos $\binom43=4$; valores de las tres sueltas $\binom{12}3=220$, palos $4^{3}=64$. Total: $13\cdot4\cdot220\cdot64=732\,160$.
  \item Probabilidades: $2\,471\,040/20\,358\,520\approx0{,}1214$ y $732\,160/20\,358\,520\approx0{,}0360$.
\end{enumerate}
\resultado la doble pareja es más probable: $0{,}1214$ frente a $0{,}0360$.
\comprobacion un programa que cuenta manos por su patrón de multiplicidades de valores, $(2,2,1,1)$ y $(3,1,1,1)$, da los mismos dos números.
\end{solucion}

\begin{solucion}{eje:3-cumple}
\datos $n$ personas además de usted; 365 días.
\planteamiento complemento: ninguna de las $n$ cumple el día de usted, cada una con probabilidad $364/365$, de forma independiente.
\pasos
\begin{enumerate}[label=\arabic*.]
  \item $\PP(A_n)=1-(364/365)^{n}$.
  \item $\PP(A_n)>1/2\iff n>\ln2/(-\ln(364/365))\approx252{,}65$. Luego $n=253$: $\PP(A_{252})\approx0{,}4991$ y $\PP(A_{253})\approx0{,}5005$.
  \item (b) En el problema clásico cuenta cualquiera de los $\binom{n}{2}$ pares; aquí solo los $n$ pares que le incluyen a usted. Por la cota de la unión, $\PP(A_n)\le n/365$, luego hace falta $n\ge183>365/2$; y $n$ debe ser aún mayor porque las otras personas coinciden a veces entre sí y cubren menos de $n$ días distintos.
\end{enumerate}
\resultado $\PP(A_n)=1-(364/365)^{n}$ y $n=253$.
\comprobacion con $n=253$ hay tantos pares que le incluyen como pares hay en un grupo de 23 personas, $\binom{23}2=253$; por eso ambos problemas tienen probabilidades cercanas a $1/2$.
\end{solucion}

\begin{solucion}{eje:3-urna}
\datos $r$ rojas, $v$ verdes, dos extracciones sin reemplazamiento.
\planteamiento resultados ordenados $(b_1,b_2)$ de bolas distintas, todos equiprobables: $(r+v)(r+v-1)$.
\pasos
\begin{enumerate}[label=\arabic*.]
  \item Primera verde: $v(r+v-1)$ resultados; probabilidad $v/(r+v)$.
  \item Segunda verde: primera roja y segunda verde, $rv$; o ambas verdes, $v(v-1)$. Total $v(r+v-1)$: misma probabilidad $v/(r+v)$.
  \item (b) Mismo color: $\frac{r(r-1)+v(v-1)}{16\cdot15}=\frac12$, luego $r^{2}+v^{2}-16=120$ y $r^{2}+v^{2}=136$.
  \item Con $v=16-r$: $r^{2}+(16-r)^{2}=136\iff r^{2}-16r+60=0\iff r\in\{6,10\}$.
\end{enumerate}
\resultado (a) ambas valen $v/(r+v)$; (b) $(r,v)=(6,10)$ o $(10,6)$.
\comprobacion con $r=6$, $v=10$: $\frac{30+90}{240}=\frac12$.
\end{solucion}

\begin{solucion}{eje:3-geometrica}
\datos $(x,y)$ uniforme en $[0,1]^{2}$.
\planteamiento calcular áreas de regiones del cuadrado.
\pasos
\begin{enumerate}[label=\arabic*.]
  \item $\{y>x+1/3\}$ es el triángulo de vértices $(0,1/3)$, $(0,1)$, $(2/3,1)$, de catetos $2/3$: área $\frac12\cdot\frac49=\frac29$. Por simetría, $\{x>y+1/3\}$ tiene la misma área. Total: $4/9$.
  \item $\{x+y<1/2\}$ es el triángulo de vértices $(0,0)$, $(1/2,0)$, $(0,1/2)$: área $\frac12\cdot\frac14=\frac18$.
\end{enumerate}
\resultado $4/9$ y $1/8$.
\comprobacion el mismo razonamiento con un umbral $d\in[0,1]$ da $\PP(\abs{x-y}>d)=(1-d)^{2}$, que vale 1 para $d=0$ y 0 para $d=1$, como debe ser; con $d=1/3$ se recupera $4/9$.
\end{solucion}

\begin{solucion}{eje:3-bonferroni}
\datos sucesos $A$, $B$, $C$.
\planteamiento \cref{prop:t3-props}(vi) y (v); después, aplicar el caso de dos sucesos dos veces.
\pasos
\begin{enumerate}[label=\arabic*.]
  \item $\PP(A\cap B)=\PP(A)+\PP(B)-\PP(A\cup B)\ge\PP(A)+\PP(B)-1$, porque $\PP(A\cup B)\le1$.
  \item Con $A\cap B$ y $C$: $\PP(A\cap B\cap C)\ge\PP(A\cap B)+\PP(C)-1$.
  \item Sustituyendo el paso 1: $\PP(A\cap B\cap C)\ge\PP(A)+\PP(B)+\PP(C)-2$.
\end{enumerate}
\resultado las dos desigualdades.
\comprobacion si $\PP(A)=\PP(B)=\PP(C)=0{,}9$, la cota garantiza $\PP(A\cap B\cap C)\ge0{,}7$; con $A$, $B$, $C$ independientes el valor sería $0{,}729$, compatible con la cota.
\end{solucion}

\begin{solucion}{eje:3-interseccion}
\datos sucesos $A_1,\dots,A_n$.
\planteamiento De Morgan y la cota de la unión (\cref{teo:t3-boole}).
\pasos
\begin{enumerate}[label=\arabic*.]
  \item $\bigl(\bigcap_iA_i\bigr)^{c}=\bigcup_iA_i^{c}$, luego $\PP\bigl(\bigcap_iA_i\bigr)=1-\PP\bigl(\bigcup_iA_i^{c}\bigr)$.
  \item Por la cota de la unión, $\PP\bigl(\bigcup_iA_i^{c}\bigr)\le\sum_i\PP(A_i^{c})$, de donde la desigualdad.
  \item Aplicación: $A_i$ = «el componente $i$ no falla», $\PP(A_i^{c})=0{,}001$. Entonces $\PP(\text{ninguno falla})\ge1-20\cdot0{,}001=0{,}98$.
\end{enumerate}
\resultado $\PP\bigl(\bigcap A_i\bigr)\ge1-\sum\PP(A_i^{c})$; sin hipótesis de independencia, la probabilidad de que no falle ninguno es al menos $0{,}98$.
\comprobacion con independencia sería $0{,}999^{20}\approx0{,}9802$, apenas por encima de la cota.
\end{solucion}

\begin{solucion}{eje:3-distancia}
\datos sucesos $A,B,C$; $d(A,B)=\PP(A\dsim B)$.
\planteamiento (a) descomponer en partes disjuntas; (b) probar $A\dsim C\subseteq(A\dsim B)\cup(B\dsim C)$ y aplicar monotonía y la cota de la unión.
\pasos
\begin{enumerate}[label=\arabic*.]
  \item (a) $A\dsim B$ es la unión disjunta de $A\setminus B$ y $B\setminus A$, y $\PP(A\setminus B)=\PP(A)-\PP(A\cap B)$ por la \cref{prop:t3-props}(iv) aplicada a $A\cap B\subseteq A$. Análogamente para $B\setminus A$. Sumando: $\PP(A)+\PP(B)-2\PP(A\cap B)$.
  \item (b) Sea $\omega\in A\setminus C$. Si $\omega\in B$, entonces $\omega\in B\setminus C\subseteq B\dsim C$; si $\omega\notin B$, entonces $\omega\in A\setminus B\subseteq A\dsim B$. El caso $\omega\in C\setminus A$ es simétrico.
  \item Así $A\dsim C\subseteq(A\dsim B)\cup(B\dsim C)$ y $d(A,C)\le\PP\bigl((A\dsim B)\cup(B\dsim C)\bigr)\le d(A,B)+d(B,C)$.
\end{enumerate}
\resultado las dos afirmaciones; $d$ se comporta como una distancia entre sucesos.
\comprobacion con dos dados, $A$ «primer dado par», $B$ «segundo dado par»: $d(A,B)=\frac12+\frac12-2\cdot\frac14=\frac12$, y en efecto $A\dsim B$ = «exactamente un dado par», de probabilidad $1/2$.
\end{solucion}

\begin{solucion}{eje:3-hash}
\datos $n$ claves, $N$ casillas, asignación uniforme e independiente.
\planteamiento \cref{prop:t3-cumple}: $\PP(\text{ninguna colisión})\le\exp(-n(n-1)/(2N))$.
\pasos
\begin{enumerate}[label=\arabic*.]
  \item (a) Si $n(n-1)\ge2N\ln2$, entonces $\exp\bigl(-\frac{n(n-1)}{2N}\bigr)\le e^{-\ln2}=\frac12$, luego $\PP(\text{alguna colisión})\ge\frac12$.
  \item (b) $2N\ln2=2\cdot2^{32}\ln2\approx5\,954\,088\,944$. Para $n=77\,163$, $n(n-1)=5\,954\,051\,406$, insuficiente; para $n=77\,164$, $n(n-1)=5\,954\,205\,732$, suficiente.
  \item $\sqrt N=65\,536$ y $77\,164/65\,536\approx1{,}177\approx\sqrt{2\ln2}$.
\end{enumerate}
\resultado $n=77\,164$ claves bastan para una probabilidad de colisión de al menos $1/2$, del orden de $\sqrt N$.
\comprobacion el producto exacto $\prod_{i<n}(1-i/N)$ con $n=77\,164$ vale $0{,}49999$, de modo que la cota es casi exacta en este régimen.
\end{solucion}

\begin{solucion}{eje:3-laplace}
\datos (a) $\Omega$ finito, $\PP(A)=\abs A/\abs\Omega$; (b) $\Omega=\N$.
\planteamiento (a) verificar los tres axiomas; (b) reducción al absurdo con la aditividad numerable.
\pasos
\begin{enumerate}[label=\arabic*.]
  \item (a) $\abs A\ge0$ da la no negatividad; $\PP(\Omega)=\abs\Omega/\abs\Omega=1$.
  \item Si $A_1,A_2,\dots$ son disjuntos dos a dos, como $\Omega$ es finito solo un número finito de ellos es no vacío; por la regla de la suma, $\abs{\bigcup_iA_i}=\sum_i\abs{A_i}$ (los vacíos aportan 0). Dividiendo por $\abs\Omega$ se obtiene la aditividad.
  \item (b) Supóngase $\PP(\{n\})=c$ para todo $n\in\N$. Por la aditividad numerable, $1=\PP(\N)=\sum_{n\in\N}c$.
  \item Si $c=0$, la serie vale $0\neq1$; si $c>0$, diverge. En ambos casos, contradicción.
\end{enumerate}
\resultado el modelo de Laplace es una ley de probabilidad, y no existe una ley «uniforme» sobre $\N$.
\comprobacion sobre $\N$ sí existen leyes que dan probabilidad positiva a todos los puntos, como $\PP(\{n\})=2^{-(n+1)}$, pero no pueden ser constantes.
\end{solucion}

% ---------------------------------------------------------------------
\section{Tema 4}

\begin{solucion}{eje:4-tabla}
\datos 100 estudiantes: 40 de Ingeniería (30 aprueban), 60 de Matemáticas (36 aprueban).
\planteamiento modelo de Laplace y definición de probabilidad condicionada.
\pasos
\begin{enumerate}[label=\arabic*.]
  \item $\PP(\text{ap})=(30+36)/100=0{,}66$.
  \item $\PP(\text{ap}\mid I)=30/40=0{,}75$.
  \item $\PP(I\mid\text{ap})=30/66=5/11\approx0{,}4545$.
  \item $\PP(\text{ap}\mid I)=0{,}75\neq0{,}66=\PP(\text{ap})$: no son independientes.
\end{enumerate}
\resultado $0{,}66$; $0{,}75$; $5/11$; no independientes.
\comprobacion Bayes: $\PP(I\mid\text{ap})=\PP(\text{ap}\mid I)\PP(I)/\PP(\text{ap})=0{,}75\cdot0{,}4/0{,}66=5/11$.
\end{solucion}

\begin{solucion}{eje:4-cachorros}
\datos dos crías, sexos equiprobables e independientes: $\{\mathrm{HH},\mathrm{HM},\mathrm{MH},\mathrm{MM}\}$ equiprobables (la primera letra, la cría mayor).
\planteamiento definición de probabilidad condicionada (\cref{ej:t4-hijos}).
\pasos
\begin{enumerate}[label=(\alph*)]
  \item Condicionante «mayor hembra» $=\{\mathrm{HH},\mathrm{HM}\}$; $\PP(\mathrm{HH}\mid\cdot)=\frac{1/4}{1/2}=\frac12$.
  \item Condicionante «al menos un macho» $=\{\mathrm{HM},\mathrm{MH},\mathrm{MM}\}$; $\PP(\mathrm{MM}\mid\cdot)=\frac{1/4}{3/4}=\frac13$.
\end{enumerate}
\resultado $1/2$ y $1/3$.
\comprobacion en (a) la información se refiere a una cría concreta y la otra sigue siendo equiprobable; en (b) no, y el suceso condicionante tiene tres casos.
\end{solucion}

\begin{solucion}{eje:4-spam}
\datos $\PP(S)=0{,}8$, $\PP(F\mid S)=0{,}1$, $\PP(F\mid S^{c})=0{,}01$.
\planteamiento teorema de Bayes con la partición $\{S,S^{c}\}$.
\pasos
\begin{enumerate}[label=\arabic*.]
  \item $\PP(F)=0{,}1\cdot0{,}8+0{,}01\cdot0{,}2=0{,}08+0{,}002=0{,}082$.
  \item $\PP(S\mid F)=0{,}08/0{,}082=40/41$.
\end{enumerate}
\resultado $40/41\approx0{,}9756$.
\comprobacion forma de razones: razón a priori $0{,}8/0{,}2=4$, razón de verosimilitudes $0{,}1/0{,}01=10$, razón a posteriori $40$, es decir, probabilidad $40/41$.
\end{solucion}

\begin{solucion}{eje:4-dos-pruebas}
\datos $\PP(E)=0{,}01$; $\PP(+\mid E)=0{,}95$; $\PP(+\mid E^{c})=0{,}05$; independencia condicional de las pruebas.
\planteamiento Bayes con la evidencia conjunta $++$, y Bayes dos veces sucesivas.
\pasos
\begin{enumerate}[label=\arabic*.]
  \item (a) $\PP(++\mid E)=0{,}95^{2}=0{,}9025$ y $\PP(++\mid E^{c})=0{,}05^{2}=0{,}0025$.
  \item $\PP(E\mid++)=\frac{0{,}01\cdot0{,}9025}{0{,}01\cdot0{,}9025+0{,}99\cdot0{,}0025}=\frac{0{,}009025}{0{,}0115}\approx0{,}7848$.
  \item (b) Tras la primera prueba: $p_1=\frac{0{,}0095}{0{,}0095+0{,}0495}=\frac{0{,}0095}{0{,}059}\approx0{,}1610$.
  \item Tomando $p_1$ como nueva probabilidad a priori: $\frac{0{,}95\,p_1}{0{,}95\,p_1+0{,}05(1-p_1)}\approx\frac{0{,}15297}{0{,}15297+0{,}04195}\approx0{,}7848$.
\end{enumerate}
\resultado $361/460\approx0{,}7848$ por ambas vías.
\comprobacion razones: a priori $1/99$; cada prueba multiplica por $0{,}95/0{,}05=19$; a posteriori $361/99$, es decir, probabilidad $361/460$.
\end{solucion}

\begin{solucion}{eje:4-piezas}
\datos tipo antiguo $V$ o nuevo $N$ con probabilidad $1/2$; $\PP(\text{def}\mid V)=0{,}15$, $\PP(\text{def}\mid N)=0{,}05$; dos extracciones con reemplazamiento, independientes dado el tipo.
\planteamiento probabilidad total y Bayes.
\pasos
\begin{enumerate}[label=\arabic*.]
  \item $\PP(DD\mid V)=0{,}15^{2}=0{,}0225$ y $\PP(DD\mid N)=0{,}05^{2}=0{,}0025$.
  \item (a) $\PP(DD)=\frac12(0{,}0225+0{,}0025)=0{,}0125$.
  \item (b) $\PP(V\mid DD)=\frac{0{,}5\cdot0{,}0225}{0{,}0125}=\frac{0{,}01125}{0{,}0125}=0{,}9$.
\end{enumerate}
\resultado $0{,}0125$ y $0{,}9$.
\comprobacion razón de verosimilitudes $0{,}0225/0{,}0025=9$ con razón a priori 1: razón a posteriori 9, probabilidad $9/10$. El número de piezas de cada tipo no interviene, porque el tipo se elige con la moneda.
\end{solucion}

\begin{solucion}{eje:4-perro}
\datos $\PP(A)=0{,}4$, $\PP(B)=0{,}6$; $\PP(\text{hallar}\mid A,\text{buscar en }A)=0{,}25$; $\PP(\text{hallar}\mid B,\text{buscar en }B)=0{,}15$.
\planteamiento regla de la multiplicación, Bayes y probabilidad total.
\pasos
\begin{enumerate}[label=\arabic*.]
  \item (a) Buscando en A: $0{,}4\cdot0{,}25=0{,}10$; en B: $0{,}6\cdot0{,}15=0{,}09$. Conviene A.
  \item (b) Sea $F$ «no se encuentra buscando en A»: $\PP(F\mid A)=0{,}75$, $\PP(F\mid B)=1$. $\PP(A\mid F)=\frac{0{,}4\cdot0{,}75}{0{,}4\cdot0{,}75+0{,}6\cdot1}=\frac{0{,}3}{0{,}9}=\frac13$.
  \item (c) Sea $H$ «se halla el primer día». $\PP(H\mid\text{buscar en A})=0{,}1$, $\PP(H\mid\text{buscar en B})=0{,}09$. $\PP(\text{buscó en A}\mid H)=\frac{0{,}5\cdot0{,}1}{0{,}5\cdot0{,}1+0{,}5\cdot0{,}09}=\frac{10}{19}$.
\end{enumerate}
\resultado (a) en A; (b) $1/3$; (c) $10/19\approx0{,}526$.
\comprobacion en (b), la búsqueda fallida en A reduce la probabilidad de A de $0{,}4$ a $0{,}33$, como cabe esperar.
\end{solucion}

\begin{solucion}{eje:4-monty7}
\datos $n$ puertas, el presentador abre $m$ con cabra entre las $n-1$ no elegidas; caso (a): $n=7$, $m=3$.
\planteamiento probabilidad total según la elección inicial sea correcta o no.
\pasos
\begin{enumerate}[label=\arabic*.]
  \item Si la elección inicial es correcta (probabilidad $1/n$), cambiar pierde siempre.
  \item Si es incorrecta (probabilidad $(n-1)/n$), el coche está tras una de las $n-1-m$ puertas cerradas restantes, y cambiando a una de ellas al azar se acierta con probabilidad $1/(n-m-1)$.
  \item $\PP(\text{ganar cambiando})=\frac{n-1}{n}\cdot\frac{1}{n-m-1}$.
  \item (a) $n=7$, $m=3$: $\frac67\cdot\frac13=\frac27$, frente a $\frac17$ sin cambiar.
  \item (b) Como $n-1>n-m-1$, $\frac{n-1}{n(n-m-1)}>\frac1n$: siempre conviene cambiar.
\end{enumerate}
\resultado $2/7$ con 7 puertas; en general $\frac{n-1}{n(n-m-1)}>\frac1n$.
\comprobacion $n=3$, $m=1$: $\frac23\cdot1=\frac23$, el valor del \cref{ej:t4-monty}.
\end{solucion}

\begin{solucion}{eje:4-monty-p}
\datos elección inicial puerta 1; $\PP(C_i)=1/3$; $\PP(M_2\mid C_1)=p$, $\PP(M_3\mid C_1)=1-p$; $\PP(M_3\mid C_2)=1$; $\PP(M_2\mid C_3)=1$.
\planteamiento Bayes con la partición $\{C_1,C_2,C_3\}$.
\pasos
\begin{enumerate}[label=\arabic*.]
  \item (a) Cambiar gana si y solo si el coche no está en la puerta~1: $2/3$.
  \item (b) $\PP(M_2)=\frac13p+0+\frac13=\frac{1+p}3$; cambiar gana si $C_3$: $\PP(C_3\mid M_2)=\frac{1/3}{(1+p)/3}=\frac1{1+p}$.
  \item (c) $\PP(M_3)=\frac13(1-p)+\frac13=\frac{2-p}3$; $\PP(C_2\mid M_3)=\frac1{2-p}$.
\end{enumerate}
\resultado $2/3$; $1/(1+p)$; $1/(2-p)$.
\comprobacion probabilidad total: $\frac{1+p}3\cdot\frac1{1+p}+\frac{2-p}3\cdot\frac1{2-p}=\frac23$. Con $p=1/2$, (b) y (c) valen $2/3$; con $p=1$, (b) vale $1/2$ y (c) vale~1.
\end{solucion}

\begin{solucion}{eje:4-paridad}
\datos dos dados; $A$, $B$, $C$ como en el enunciado.
\planteamiento la suma es impar si y solo si exactamente uno de los dos es impar.
\pasos
\begin{enumerate}[label=\arabic*.]
  \item $\PP(A)=\PP(B)=\frac12$; $\PP(C)=\PP(\text{impar, par})+\PP(\text{par, impar})=\frac14+\frac14=\frac12$.
  \item $\PP(A\cap B)=\frac14$; $A\cap C$ = «primero impar, segundo par», $\frac14$; $B\cap C$ = «primero par, segundo impar», $\frac14$. Independientes dos a dos.
  \item $A\cap B\cap C=\emptyset$: si ambos son impares, la suma es par. $\PP(A\cap B\cap C)=0\neq\frac18$.
\end{enumerate}
\resultado independientes dos a dos, pero no mutuamente.
\comprobacion es el mismo fenómeno que el \cref{ej:t4-dosados}: $C$ es una función de las paridades de los dos dados.
\end{solucion}

\begin{solucion}{eje:4-ajedrez}
\datos niveles $P$, $I$, $M$ con probabilidad $1/3$; $\PP(G\mid P)=0{,}9$, $\PP(G\mid I)=0{,}5$, $\PP(G\mid M)=0{,}3$; partidas independientes dado el nivel.
\planteamiento probabilidad total, condicionando al nivel.
\pasos
\begin{enumerate}[label=\arabic*.]
  \item (a) $\PP(G_1)=\frac13(0{,}9+0{,}5+0{,}3)=\frac{1{,}7}3\approx0{,}5667$.
  \item $\PP(G_1\cap G_2)=\frac13(0{,}9^{2}+0{,}5^{2}+0{,}3^{2})=\frac{1{,}15}3\approx0{,}3833$.
  \item (b) $\PP(G_2\mid G_1)=\frac{1{,}15}{1{,}7}=\frac{23}{34}\approx0{,}6765$.
  \item (c) $\PP(G_2\mid G_1)\neq\PP(G_2)=0{,}5667$: no son independientes. Ganar la primera hace más probable que el rival sea débil y, por tanto, que se gane la segunda. Son independientes \emph{condicionalmente} al nivel, pero no incondicionalmente.
\end{enumerate}
\resultado $17/30$; $23/34$; dependientes, aunque condicionalmente independientes dado el nivel.
\comprobacion $\PP(G_1\cap G_2)=0{,}3833>\PP(G_1)^{2}=0{,}3211$: los resultados están positivamente relacionados.
\end{solucion}

\begin{solucion}{eje:4-si-mismo}
\datos un suceso $A$.
\planteamiento la independencia de $A$ consigo mismo es $\PP(A\cap A)=\PP(A)^{2}$.
\pasos
\begin{enumerate}[label=\arabic*.]
  \item $\PP(A)=\PP(A)^{2}\iff\PP(A)(1-\PP(A))=0\iff\PP(A)\in\{0,1\}$.
  \item Si $\PP(A)=0$: $0\le\PP(A\cap B)\le\PP(A)=0=\PP(A)\PP(B)$.
  \item Si $\PP(A)=1$: $\PP(A^{c})=0$, luego $\PP(A^{c}\cap B)=0$ y $\PP(A\cap B)=\PP(B)-\PP(A^{c}\cap B)=\PP(B)=\PP(A)\PP(B)$.
\end{enumerate}
\resultado $A$ es independiente de sí mismo si y solo si $\PP(A)\in\{0,1\}$, y entonces es independiente de cualquier suceso.
\comprobacion $\Omega$ y $\emptyset$ son independientes de todo suceso, como cabe esperar: saber que ha ocurrido $\Omega$ no informa de nada.
\end{solucion}

\begin{solucion}{eje:4-union-indep}
\datos (a) $A$, $B$, $C$ mutuamente independientes; (b) $A_1,\dots,A_n$ independientes.
\planteamiento (a) distributiva e inclusión--exclusión; (b) independencia de los complementarios y De Morgan.
\pasos
\begin{enumerate}[label=\arabic*.]
  \item (a) $\PP(A\cap(B\cup C))=\PP((A\cap B)\cup(A\cap C))=\PP(A\cap B)+\PP(A\cap C)-\PP(A\cap B\cap C)$.
  \item Por la independencia mutua, esto es $\PP(A)\bigl[\PP(B)+\PP(C)-\PP(B)\PP(C)\bigr]=\PP(A)\PP(B\cup C)$, usando $\PP(B\cap C)=\PP(B)\PP(C)$.
  \item (b) Si $A_1,\dots,A_n$ son independientes, sustituir uno de ellos, $A_i$, por $A_i^{c}$ conserva la independencia: para todo $S\ni i$, con $T=S\setminus\{i\}$, $\PP\bigl(\bigcap_{T}A_j\cap A_i^{c}\bigr)=\PP\bigl(\bigcap_TA_j\bigr)-\PP\bigl(\bigcap_SA_j\bigr)=\prod_T\PP(A_j)\,(1-\PP(A_i))$. Repitiendo, $A_1^{c},\dots,A_n^{c}$ son independientes.
  \item $\PP\bigl(\bigcup A_i\bigr)=1-\PP\bigl(\bigcap A_i^{c}\bigr)=1-\prod_i(1-\PP(A_i))$.
\end{enumerate}
\resultado las dos afirmaciones.
\comprobacion $n$ lanzamientos de una moneda, $A_i$ «cara en el $i$»: $\PP(\text{alguna cara})=1-2^{-n}$, el resultado conocido.
\end{solucion}

\begin{solucion}{eje:4-producto-triple}
\planteamiento basta que la intersección triple tenga la probabilidad «correcta» mientras falla alguna de las condiciones dos a dos.
\pasos
\begin{enumerate}[label=\arabic*.]
  \item Ejemplo 1: $\Omega=[8]$ con ley uniforme, $A=\{1,2,3,4\}$, $B=\{1,2,3,5\}$, $C=\{1,6,7,8\}$. Cada uno tiene probabilidad $1/2$.
  \item $A\cap B\cap C=\{1\}$, con probabilidad $1/8=\PP(A)\PP(B)\PP(C)$.
  \item Pero $A\cap B=\{1,2,3\}$ tiene probabilidad $3/8\neq1/4$: no son mutuamente independientes (ni siquiera dos a dos).
  \item Ejemplo 2, degenerado: $C=\emptyset$ y $A=B$ con $\PP(A)=1/2$. La triple vale $0=\PP(A)\PP(B)\cdot0$, pero $\PP(A\cap B)=1/2\neq1/4$.
\end{enumerate}
\resultado la condición sobre la intersección triple no implica la independencia mutua.
\comprobacion en el ejemplo 1, $\PP(A\cap C)=\PP(\{1\})=1/8\neq1/4$ y $\PP(B\cap C)=1/8\neq1/4$: fallan las tres condiciones dobles.
\end{solucion}

\begin{solucion}{eje:4-razones}
\datos sucesos $B$, $E_1$, $E_2$ con probabilidades positivas; en (b), independencia condicional dada $B$ y dada $B^{c}$.
\planteamiento (a) cociente de dos fórmulas de Bayes; (b) comparar el producto de las dos razones de verosimilitudes con la conjunta.
\pasos
\begin{enumerate}[label=\arabic*.]
  \item (a) $\PP(B\mid A)=\PP(A\mid B)\PP(B)/\PP(A)$ y $\PP(B^{c}\mid A)=\PP(A\mid B^{c})\PP(B^{c})/\PP(A)$; al dividir se cancela $\PP(A)$.
  \item (b) Una sola actualización: $O_{12}=O_0\cdot\frac{\PP(E_1\cap E_2\mid B)}{\PP(E_1\cap E_2\mid B^{c})}=O_0\cdot\frac{\PP(E_1\mid B)\PP(E_2\mid B)}{\PP(E_1\mid B^{c})\PP(E_2\mid B^{c})}$, por la independencia condicional.
  \item Dos actualizaciones: tras $E_1$, $O_1=O_0\frac{\PP(E_1\mid B)}{\PP(E_1\mid B^{c})}$. Al incorporar $E_2$, las verosimilitudes son ahora $\PP(E_2\mid B\cap E_1)$ y $\PP(E_2\mid B^{c}\cap E_1)$, que por la independencia condicional valen $\PP(E_2\mid B)$ y $\PP(E_2\mid B^{c})$.
  \item Por tanto $O_2=O_1\frac{\PP(E_2\mid B)}{\PP(E_2\mid B^{c})}=O_{12}$.
\end{enumerate}
\resultado ambas formas de actualizar coinciden.
\comprobacion el \cref{eje:4-dos-pruebas} es un caso numérico: $\frac1{99}\cdot19\cdot19=\frac{361}{99}$ por ambas vías.
\end{solucion}

\begin{solucion}{eje:4-simetria}
\datos $0<\PP(A),\PP(B)<1$.
\planteamiento reescribir cada afirmación en términos de $\PP(A\cap B)$ y $\PP(A)\PP(B)$.
\pasos
\begin{enumerate}[label=\arabic*.]
  \item (a) $\iff\PP(A\cap B)/\PP(B)>\PP(A)\iff\PP(A\cap B)>\PP(A)\PP(B)$.
  \item (b) $\iff\PP(A\cap B)/\PP(A)>\PP(B)\iff\PP(A\cap B)>\PP(A)\PP(B)$. Luego (a)$\iff$(b).
  \item Por la probabilidad total, $\PP(A)=\PP(A\mid B)\PP(B)+\PP(A\mid B^{c})(1-\PP(B))$, de donde
  \[
  \PP(A\mid B)-\PP(A)=(1-\PP(B))\bigl(\PP(A\mid B)-\PP(A\mid B^{c})\bigr).
  \]
  \item Como $1-\PP(B)>0$, (a)$\iff$(c).
\end{enumerate}
\resultado las tres afirmaciones son equivalentes: la «correlación positiva» entre sucesos es simétrica.
\comprobacion en el \cref{eje:4-tabla}: $\PP(\text{ap}\mid I)=0{,}75>0{,}66$ y $\PP(I\mid\text{ap})=5/11\approx0{,}45>0{,}4=\PP(I)$; además $\PP(\text{ap}\mid I^{c})=0{,}6<0{,}75$.
\end{solucion}

% ---------------------------------------------------------------------
\section{Tema 5}

\begin{solucion}{eje:5-masa}
\datos $\PP(X=k)=ck$, $k=1,\dots,4$.
\planteamiento normalización, acumulación y definiciones de esperanza y varianza (\cref{prop:t5-var}).
\pasos
\begin{enumerate}[label=\arabic*.]
  \item $c(1+2+3+4)=10c=1$, luego $c=1/10$.
  \item $F_X(1)=0{,}1$, $F_X(2)=0{,}3$, $F_X(3)=0{,}6$, $F_X(4)=1$.
  \item $\E[X]=\sum k\cdot\frac{k}{10}=\frac{1+4+9+16}{10}=3$.
  \item $\E[X^{2}]=\sum k^{2}\cdot\frac k{10}=\frac{1+8+27+64}{10}=10$; $\Var(X)=10-9=1$.
\end{enumerate}
\resultado $c=1/10$; $F_X=0{,}1;\,0{,}3;\,0{,}6;\,1$; $\E[X]=3$; $\Var(X)=1$.
\comprobacion $\Var(X)=\sum(k-3)^{2}\frac k{10}=\frac{4+2+0+4}{10}=1$.
\end{solucion}

\begin{solucion}{eje:5-diferencia}
\datos dos dados de tres caras; $X=D_1-D_2$; 9 resultados equiprobables.
\planteamiento contar los pares que dan cada diferencia; LOTUS para $X^{2}$.
\pasos
\begin{enumerate}[label=\arabic*.]
  \item Valores: $-2$ (1 par), $-1$ (2), $0$ (3), $1$ (2), $2$ (1); $p_X=\frac19,\frac29,\frac39,\frac29,\frac19$.
  \item $\E[X]=0$ por simetría ($X$ y $-X$ tienen la misma distribución).
  \item $\Var(X)=\E[X^{2}]=\frac{4\cdot1+1\cdot2+0+1\cdot2+4\cdot1}{9}=\frac{12}9=\frac43$.
  \item $X^{2}$ toma $0$ con probabilidad $\frac39$, $1$ con $\frac49$ y $4$ con $\frac29$; $\E[X^{2}]=\frac49+\frac89=\frac43$.
\end{enumerate}
\resultado $p_X$ como arriba; $\E[X]=0$; $\Var(X)=4/3$; $p_{X^{2}}=(\frac13,\frac49,\frac29)$ en $(0,1,4)$.
\comprobacion por independencia, $\Var(D_1-D_2)=\Var(D_1)+\Var(D_2)=2\cdot\frac{3^{2}-1}{12}=\frac43$ (\cref{tab:t5-discretas}).
\end{solucion}

\begin{solucion}{eje:5-binomial}
\datos $X\sim\Bin(10;0{,}3)$.
\planteamiento \cref{teo:t5-binomial}.
\pasos
\begin{enumerate}[label=\arabic*.]
  \item $\PP(X=3)=\binom{10}3\,0{,}3^{3}\,0{,}7^{7}=120\cdot0{,}027\cdot0{,}0823543\approx0{,}2668$.
  \item $\PP(X\ge1)=1-0{,}7^{10}\approx1-0{,}0282=0{,}9718$.
  \item $\E[X]=10\cdot0{,}3=3$; $\Var(X)=10\cdot0{,}3\cdot0{,}7=2{,}1$.
\end{enumerate}
\resultado $0{,}2668$; $0{,}9718$; $3$; $2{,}1$.
\comprobacion la moda de la binomial es $\lfloor(n+1)p\rfloor=3$, coherente con que $\PP(X=3)$ sea el valor más alto de la función de masa.
\end{solucion}

\begin{solucion}{eje:5-geometrica}
\datos $X\sim\Geom(0{,}2)$, $q=0{,}8$.
\planteamiento \cref{prop:t5-geom}.
\pasos
\begin{enumerate}[label=\arabic*.]
  \item $\PP(X=3)=q^{2}p=0{,}64\cdot0{,}2=0{,}128$.
  \item $\PP(X>5)=q^{5}=0{,}32768$.
  \item $\E[X]=1/0{,}2=5$; $\Var(X)=0{,}8/0{,}04=20$.
  \item Falta de memoria: $\PP(X>8\mid X>3)=\PP(X>5)=0{,}32768$.
\end{enumerate}
\resultado $0{,}128$; $0{,}32768$; $5$; $20$; $0{,}32768$.
\comprobacion directamente, $\PP(X>8)/\PP(X>3)=0{,}8^{8}/0{,}8^{3}=0{,}8^{5}$.
\end{solucion}

\begin{solucion}{eje:5-poisson-aprox}
\datos $X\sim\Bin(1000;\,0{,}002)$; aproximación $\Poi(2)$.
\planteamiento \cref{teo:t5-poisson} con $\lambda=np=2$.
\pasos
\begin{enumerate}[label=\arabic*.]
  \item Binomial: $\PP(X=0)=0{,}998^{1000}\approx0{,}13506$; $\PP(X=1)=1000\cdot0{,}002\cdot0{,}998^{999}\approx0{,}27067$; $\PP(X\ge2)\approx1-0{,}13506-0{,}27067=0{,}59427$.
  \item Poisson: $\PP(X=0)=e^{-2}\approx0{,}13534$; $\PP(X=1)=2e^{-2}\approx0{,}27067$; $\PP(X\ge2)=1-3e^{-2}\approx0{,}59399$.
\end{enumerate}
\resultado la aproximación de Poisson acierta en las tres primeras cifras decimales.
\comprobacion las dos distribuciones tienen media 2 y varianzas $1{,}996$ y $2$, casi iguales.
\end{solucion}

\begin{solucion}{eje:5-cumple-esp}
\datos 50 personas, 365 días.
\planteamiento linealidad con indicadoras (\cref{prop:t5-puente}): una por pareja en (a) y una por día en (b).
\pasos
\begin{enumerate}[label=\arabic*.]
  \item (a) Para cada una de las $\binom{50}2=1225$ parejas, $\PP(\text{mismo día})=1/365$. $\E=1225/365\approx3{,}356$.
  \item (b) Para el día $d$, sea $N_d\sim\Bin(50,1/365)$ el número de personas nacidas ese día. $\PP(N_d\ge2)=1-\bigl(\tfrac{364}{365}\bigr)^{50}-50\cdot\tfrac1{365}\bigl(\tfrac{364}{365}\bigr)^{49}$.
  \item $\E=365\,\PP(N_d\ge2)\approx3{,}076$.
\end{enumerate}
\resultado $\approx3{,}356$ parejas y $\approx3{,}076$ días.
\comprobacion (b) es algo menor que (a) porque un día con tres personas aporta 3 parejas pero un solo día.
\end{solucion}

\begin{solucion}{eje:5-hash}
\datos $k$ claves, $n$ casillas.
\planteamiento indicadoras por casilla; el número de claves en una casilla es $\Bin(k,1/n)$.
\pasos
\begin{enumerate}[label=\arabic*.]
  \item Vacías: $\PP(\text{casilla vacía})=(1-1/n)^{k}$; esperanza $n(1-1/n)^{k}$.
  \item Con exactamente una: $\PP=k\cdot\frac1n(1-\frac1n)^{k-1}$; esperanza $k(1-1/n)^{k-1}$.
  \item Con más de una: $n-n(1-\frac1n)^{k}-k(1-\frac1n)^{k-1}$, por linealidad, pues las tres indicadoras de cada casilla suman~1.
  \item $k=n=100$: $100\cdot0{,}99^{100}\approx36{,}60$; $100\cdot0{,}99^{99}\approx36{,}97$; $100-36{,}60-36{,}97\approx26{,}42$.
\end{enumerate}
\resultado $36{,}60$ vacías, $36{,}97$ con una clave y $26{,}42$ con más de una.
\comprobacion la fracción de casillas vacías, $0{,}366$, está próxima a $e^{-1}\approx0{,}368$, el límite de $(1-1/n)^{n}$.
\end{solucion}

\begin{solucion}{eje:5-hijos}
\datos sexos equiprobables e independientes.
\planteamiento el primer hijo fija un sexo; a partir de ahí se espera el primer hijo del sexo opuesto, lo que es una geométrica.
\pasos
\begin{enumerate}[label=\arabic*.]
  \item Sea $Y$ el número de hijos tras el primero hasta el primero del sexo contrario, incluido. $Y\sim\Geom(1/2)$, independientemente del sexo del primero.
  \item $X=1+Y$, $\E[X]=1+2=3$.
\end{enumerate}
\resultado 3 hijos en promedio.
\comprobacion directamente: $\PP(X=k)=2\cdot(1/2)^{k}$ para $k\ge2$ (los $k-1$ primeros del mismo sexo, cualquiera de los dos, y el $k$-ésimo distinto), y $\sum_{k\ge2}k\,2^{1-k}=3$.
\end{solucion}

\begin{solucion}{eje:5-inversiones}
\datos permutación uniforme de $n$ números distintos.
\planteamiento una indicadora por cada par de posiciones $i<j$.
\pasos
\begin{enumerate}[label=\arabic*.]
  \item Para $i<j$, intercambiar los elementos de las posiciones $i$ y $j$ es una biyección entre las permutaciones con el par invertido y las que no lo tienen; luego $\PP(\text{invertido})=1/2$.
  \item Por linealidad, $\E[\text{inversiones}]=\binom n2\cdot\frac12=\frac{n(n-1)}4$.
\end{enumerate}
\resultado $n(n-1)/4$.
\comprobacion para $n=4$, las 24 permutaciones suman 72 inversiones, de media $3=4\cdot3/4$.
\end{solucion}

\begin{solucion}{eje:5-categorica}
\datos 10 tokens i.i.d. con $p=(0{,}5;\,0{,}3;\,0{,}2)$.
\planteamiento el número de apariciones de cada token es binomial; para (b), una indicadora por token.
\pasos
\begin{enumerate}[label=\arabic*.]
  \item $N_{\texttt a}\sim\Bin(10;0{,}5)$, $N_{\texttt b}\sim\Bin(10;0{,}3)$, $N_{\texttt c}\sim\Bin(10;0{,}2)$: medias $5$, $3$ y $2$.
  \item $\Var(N_{\texttt a})=10\cdot0{,}5\cdot0{,}5=2{,}5$.
  \item (b) $\PP(\text{el token }i\text{ aparece})=1-(1-p_i)^{10}$. Distintos esperados: $3-(0{,}5^{10}+0{,}7^{10}+0{,}8^{10})\approx3-(0{,}00098+0{,}02825+0{,}10737)=2{,}8634$.
\end{enumerate}
\resultado medias $5$, $3$, $2$; varianza $2{,}5$; $\approx2{,}863$ tokens distintos.
\comprobacion las medias suman 10, el número de tokens generados, como exige la linealidad.
\end{solucion}

\begin{solucion}{eje:5-passk}
\datos $n=10$ soluciones, $c=3$ correctas, se eligen $k$ sin reemplazamiento.
\planteamiento modelo de Laplace sobre los $\binom nk$ subconjuntos; complemento.
\pasos
\begin{enumerate}[label=\arabic*.]
  \item (a) Ninguna correcta equivale a elegir las $k$ entre las $n-c$ incorrectas: $\binom{n-c}k$ subconjuntos. $\PP(\text{alguna correcta})=1-\binom{n-c}k\big/\binom nk$.
  \item (b) $k=1$: $1-7/10=0{,}3$. $k=2$: $1-21/45=8/15\approx0{,}533$. $k=5$: $1-21/252=11/12\approx0{,}917$.
\end{enumerate}
\resultado $0{,}3$; $0{,}533$; $0{,}917$.
\comprobacion con $k=1$ se recupera la proporción de correctas, $c/n$; y si $k>n-c=7$ la probabilidad es 1, porque $\binom7k=0$.
\end{solucion}

\begin{solucion}{eje:5-colas}
\datos $X$ con valores en $\{0,1,2,\dots\}$.
\planteamiento intercambiar el orden en una suma doble de términos no negativos.
\pasos
\begin{enumerate}[label=\arabic*.]
  \item $\sum_{k\ge1}\PP(X\ge k)=\sum_{k\ge1}\sum_{j\ge k}p_X(j)$.
  \item Los pares $(k,j)$ con $1\le k\le j$ pueden sumarse primero en $k$: para cada $j\ge1$ hay $j$ valores de $k$. Como todos los términos son no negativos, el reordenamiento es lícito: $\sum_{j\ge1}j\,p_X(j)=\E[X]$.
  \item Equivalentemente, $X=\sum_{k\ge1}I_{\{X\ge k\}}$ (si $X=j$, valen 1 exactamente las indicadoras con $k\le j$).
  \item Geométrica: $\PP(X\ge k)=q^{k-1}$, luego $\E[X]=\sum_{k\ge1}q^{k-1}=\frac1{1-q}=\frac1p$.
\end{enumerate}
\resultado $\E[X]=\sum_{k\ge1}\PP(X\ge k)$ y, para la geométrica, $1/p$.
\comprobacion para un dado, $\sum_{k=1}^{6}\frac{7-k}{6}=\frac{21}6=3{,}5$.
\end{solucion}

\begin{solucion}{eje:5-poisson}
\datos $p(k)=e^{-\lambda}\lambda^{k}/k!$.
\planteamiento serie de la exponencial, $\sum_{k\ge0}\lambda^{k}/k!=e^{\lambda}$, y desplazamiento de índices.
\pasos
\begin{enumerate}[label=\arabic*.]
  \item $\sum_kp(k)=e^{-\lambda}e^{\lambda}=1$.
  \item $\E[X]=\sum_{k\ge1}k\,e^{-\lambda}\frac{\lambda^{k}}{k!}=\lambda\sum_{k\ge1}e^{-\lambda}\frac{\lambda^{k-1}}{(k-1)!}=\lambda$.
  \item $\E[X(X-1)]=\sum_{k\ge2}k(k-1)e^{-\lambda}\frac{\lambda^{k}}{k!}=\lambda^{2}\sum_{k\ge2}e^{-\lambda}\frac{\lambda^{k-2}}{(k-2)!}=\lambda^{2}$.
  \item $\Var(X)=\E[X(X-1)]+\E[X]-\E[X]^{2}=\lambda^{2}+\lambda-\lambda^{2}=\lambda$.
\end{enumerate}
\resultado la función de masa suma 1 y $\E[X]=\Var(X)=\lambda$.
\comprobacion la binomial $\Bin(n,\lambda/n)$ tiene media $\lambda$ y varianza $\lambda(1-\lambda/n)\to\lambda$, coherente con el \cref{teo:t5-poisson}.
\end{solucion}

\begin{solucion}{eje:5-sin-memoria}
\datos $X\in\{1,2,\dots\}$ sin memoria; $q=\PP(X>1)\in(0,1)$; $g(k)=\PP(X>k)$.
\planteamiento la falta de memoria se escribe $g(m+k)=g(m)g(k)$; se resuelve por inducción.
\pasos
\begin{enumerate}[label=\arabic*.]
  \item $g(0)=1$ porque $X\ge1$. Si $g(m)>0$, $\PP(X>m+k\mid X>m)=g(m+k)/g(m)$, luego $g(m+k)=g(m)g(k)$.
  \item Inducción: $g(k)=q^{k}$. Cierto para $k=0,1$. Si $g(k)=q^{k}>0$, la relación con $m=k$ y $1$ da $g(k+1)=g(k)g(1)=q^{k+1}$.
  \item $\PP(X=k)=g(k-1)-g(k)=q^{k-1}(1-q)$ para $k\ge1$.
\end{enumerate}
\resultado $X\sim\Geom(p)$ con $p=1-q$.
\comprobacion la recíproca es la \cref{prop:t5-geom}(iii): la geométrica cumple la propiedad.
\end{solucion}

\begin{solucion}{eje:5-geom-cond}
\datos $X$, $Y$ independientes $\Geom(p)$, $q=1-p$, $n\ge2$.
\planteamiento probabilidad condicionada e independencia.
\pasos
\begin{enumerate}[label=\arabic*.]
  \item Para $1\le i\le n-1$: $\PP(X=i,\,X+Y=n)=\PP(X=i)\PP(Y=n-i)=q^{i-1}p\cdot q^{n-i-1}p=p^{2}q^{n-2}$.
  \item No depende de $i$. Sumando en $i$: $\PP(X+Y=n)=(n-1)p^{2}q^{n-2}$.
  \item $\PP(X=i\mid X+Y=n)=\frac{p^{2}q^{n-2}}{(n-1)p^{2}q^{n-2}}=\frac1{n-1}$.
\end{enumerate}
\resultado uniforme en $\{1,\dots,n-1\}$.
\comprobacion interpretación: dado que el segundo éxito de una sucesión de ensayos ocurre en el ensayo $n$, el primero está en cualquiera de las $n-1$ posiciones anteriores con igual probabilidad, porque todas las sucesiones con dos éxitos en $n$ ensayos y el último éxito en $n$ tienen la misma probabilidad $p^{2}q^{n-2}$.
\end{solucion}

% ---------------------------------------------------------------------
\section{Tema 6}

\begin{solucion}{eje:6-conceptos}
\pasos
\begin{enumerate}[label=(\alph*)]
  \item Falsa: la densidad de $\Unif(0,\tfrac12)$ vale 2. Lo que debe ser $\le1$ es la integral sobre cualquier intervalo.
  \item Verdadera: $\PP(X\le2)=\PP(X<2)+\PP(X=2)$ y $\PP(X=2)=0$.
  \item Falsa: si $X$ es continua, $\{X=x\}$ tiene probabilidad 0 para todo $x$, y sin embargo $X$ toma algún valor.
  \item Verdadera: si $x\le y$, $\{X\le x\}\subseteq\{X\le y\}$ y se aplica la monotonía.
\end{enumerate}
\resultado (a) F; (b) V; (c) F; (d) V.
\comprobacion en (c), si todos los sucesos de probabilidad nula fueran imposibles, $\Omega=\bigcup_x\{X=x\}$ sería imposible; no hay contradicción con la aditividad porque la unión no es numerable.
\end{solucion}

\begin{solucion}{eje:6-densidad}
\datos $f(x)=cx^{2}$ en $[0,2]$.
\planteamiento normalización e integrales de potencias.
\pasos
\begin{enumerate}[label=\arabic*.]
  \item $\int_0^2cx^{2}dx=\frac{8c}3=1$, luego $c=3/8$.
  \item $F(x)=0$ si $x<0$; $F(x)=x^{3}/8$ si $0\le x\le2$; $F(x)=1$ si $x>2$.
  \item $\PP(X>1)=1-F(1)=7/8$.
  \item $\E[X]=\frac38\int_0^2x^{3}dx=\frac38\cdot4=\frac32$; $\E[X^{2}]=\frac38\int_0^2x^{4}dx=\frac38\cdot\frac{32}5=\frac{12}5$; $\Var(X)=\frac{12}5-\frac94=\frac3{20}$.
\end{enumerate}
\resultado $c=3/8$; $F(x)=x^{3}/8$ en $[0,2]$; $7/8$; $3/2$; $3/20$.
\comprobacion $F(2)=1$, y la media está a la derecha del punto medio, como corresponde a una densidad creciente.
\end{solucion}

\begin{solucion}{eje:6-pasas}
\datos densidad de pasas $40-h$ en $[0,30]$.
\planteamiento total por integración; la densidad de $H$ es la densidad de pasas normalizada.
\pasos
\begin{enumerate}[label=\arabic*.]
  \item (a) $\int_0^{30}(40-h)\,dh=1200-450=750$ pasas.
  \item (b) $g(h)=\frac{40-h}{750}$ en $[0,30]$; $G(h)=\frac{40h-h^{2}/2}{750}$ en $[0,30]$, 0 antes y 1 después.
  \item (c) $G(10)=\frac{400-50}{750}=\frac7{15}\approx0{,}467$.
\end{enumerate}
\resultado 750 pasas; $g$ y $G$ como arriba; $7/15$.
\comprobacion $G(30)=\frac{1200-450}{750}=1$; y $7/15>1/3$ porque las pasas se acumulan en el fondo.
\end{solucion}

\begin{solucion}{eje:6-exponencial}
\datos $f(y)=\lambda e^{-\lambda y}$, $y\ge0$.
\planteamiento integrar la densidad.
\pasos
\begin{enumerate}[label=\arabic*.]
  \item $F(y)=\int_0^y\lambda e^{-\lambda t}dt=\bigl[-e^{-\lambda t}\bigr]_0^y=1-e^{-\lambda y}$ para $y\ge0$.
  \item $\PP(Y\le1/\lambda)=1-e^{-1}\approx0{,}632$.
  \item Mediana: $1-e^{-\lambda m}=\frac12\iff m=\frac{\ln2}\lambda$.
\end{enumerate}
\resultado $F(y)=1-e^{-\lambda y}$; $1-e^{-1}$; $\ln2/\lambda$.
\comprobacion la mediana, $0{,}693/\lambda$, es menor que la media, $1/\lambda$, como corresponde a una distribución con cola larga a la derecha.
\end{solucion}

\begin{solucion}{eje:6-autobus}
\datos A: $W\sim\Unif(0,10)$; B: $W\sim\Expo(1/10)$.
\planteamiento \cref{prop:t6-unif} y \cref{prop:t6-memoria}; probabilidad condicionada.
\pasos
\begin{enumerate}[label=\arabic*.]
  \item (a) A: $\E[W]=5$; B: $\E[W]=10$.
  \item (b) A: $\PP(W\ge9\mid W\ge6)=\frac{1/10}{4/10}=\frac14$.
  \item B: $\PP(W\ge9\mid W\ge6)=\PP(W\ge3)=e^{-3/10}\approx0{,}741$.
\end{enumerate}
\resultado medias 5 y 10 minutos; $1/4$ en A y $0{,}741$ en B.
\comprobacion en A, tras 6 minutos la espera restante es uniforme en $[0,4]$ y $\PP(\ge3)=1/4$; en B, la espera restante vuelve a ser exponencial de media 10, por la falta de memoria.
\end{solucion}

\begin{solucion}{eje:6-tipificar}
\datos $X\sim\Nor(2,9)$, $\mu=2$, $\sigma=3$.
\planteamiento \cref{teo:t6-tipificacion}: $Z=(X-2)/3\sim\Nor(0,1)$.
\pasos
\begin{enumerate}[label=\arabic*.]
  \item $\PP(X\le2)=\Phi(0)=0{,}5$.
  \item $\PP(X-\mu>1{,}5\sigma)=\PP(Z>1{,}5)=1-0{,}9332=0{,}0668$.
  \item $\PP(\abs{X-\mu}<1{,}5\sigma)=2\Phi(1{,}5)-1=0{,}8664$.
\end{enumerate}
\resultado $0{,}5$; $0{,}0668$; $0{,}8664$.
\comprobacion los resultados no dependen de $\mu$ ni de $\sigma$: solo de cuántas desviaciones típicas se alejan los umbrales de la media.
\end{solucion}

\begin{solucion}{eje:6-gestacion}
\datos $X\sim\Nor(280;\,8{,}5^{2})$.
\planteamiento tipificar y usar $\Phi$ (tabla o software) y su inversa.
\pasos
\begin{enumerate}[label=\arabic*.]
  \item (a) $z=(273-280)/8{,}5\approx-0{,}824$; $\PP(X\le273)=\Phi(-0{,}824)=1-\Phi(0{,}824)\approx0{,}205$.
  \item (b) Se busca $d$ con $\PP(X\le d)=0{,}05$, es decir, $\frac{d-280}{8{,}5}=-1{,}645$ (\cref{tab:t6-phi}); $d=280-1{,}645\cdot8{,}5\approx266{,}0$.
\end{enumerate}
\resultado (a) $\approx0{,}205$; (b) $d\approx266$, unos 14 días antes de la fecha prevista.
\comprobacion (a) está entre $\Phi(-1)=0{,}159$ y $\Phi(-0{,}5)=0{,}309$, como corresponde a $z=-0{,}82$.
\end{solucion}

\begin{solucion}{eje:6-palo}
\datos $U\sim\Unif(0,1)$; $L=\max(U,1-U)\in[\tfrac12,1]$.
\planteamiento calcular $F_L(x)=\PP(L\le x)$ para $x\in[\tfrac12,1]$.
\pasos
\begin{enumerate}[label=\arabic*.]
  \item $L\le x\iff U\le x$ y $1-U\le x\iff1-x\le U\le x$.
  \item $F_L(x)=x-(1-x)=2x-1$ para $\frac12\le x\le1$ (0 antes y 1 después).
  \item Es la función de distribución de $\Unif(\tfrac12,1)$; $\E[L]=\frac34$.
\end{enumerate}
\resultado $L\sim\Unif(\tfrac12,1)$ y $\E[L]=3/4$.
\comprobacion el trozo corto, $1-L$, es $\Unif(0,\tfrac12)$ con media $1/4$, y $\frac34+\frac14=1$.
\end{solucion}

\begin{solucion}{eje:6-markov}
\datos $X\sim\Expo(1)$: $\E[X]=1$, $\Var(X)=1$.
\planteamiento \cref{teo:t6-markov,teo:t6-chebyshev}.
\pasos
\begin{enumerate}[label=\arabic*.]
  \item (a) Markov: $\PP(X\ge3)\le1/3$. Exacto: $e^{-3}\approx0{,}0498$.
  \item (b) Chebyshev: $\PP(\abs{X-1}\ge2)\le1/4$. Como $X\ge0$, $\abs{X-1}\ge2\iff X\ge3$; exacto $e^{-3}\approx0{,}0498$.
\end{enumerate}
\resultado cotas $1/3$ y $1/4$ frente al valor exacto $0{,}0498$.
\comprobacion ambas cotas son válidas pero poco ajustadas: solo usan la media o la varianza, no la forma exponencial de la cola.
\end{solucion}

\begin{solucion}{eje:6-n-cheb}
\datos media $\mu$, varianza $\sigma^{2}$, $\Var(\Xbar_n)=\sigma^{2}/n$.
\planteamiento Chebyshev para $\Xbar_n$ con $c=2\sigma$.
\pasos
\begin{enumerate}[label=\arabic*.]
  \item $\PP(\abs{\Xbar_n-\mu}\ge2\sigma)\le\frac{\sigma^{2}/n}{4\sigma^{2}}=\frac1{4n}$.
  \item $\frac1{4n}\le0{,}01\iff n\ge25$.
\end{enumerate}
\resultado $n=25$ basta, sea cual sea la distribución.
\comprobacion con $n=25$, $\sigma_{\Xbar}=\sigma/5$ y el umbral $2\sigma$ son 10 desviaciones típicas de la media muestral; Chebyshev da $1/100$.
\end{solucion}

\begin{solucion}{eje:6-unif-lgn}
\datos $U_i\sim\Unif(-1,1)$ independientes.
\planteamiento \cref{prop:t6-unif} y \cref{teo:t6-lgn}.
\pasos
\begin{enumerate}[label=\arabic*.]
  \item $\mu=0$ y $\sigma^{2}=\frac{2^{2}}{12}=\frac13$, finitas; por la ley débil, $\overline U_n\to0$ en probabilidad.
  \item $\PP(\abs{\overline U_{100}}\ge0{,}1)\le\frac{1/3}{100\cdot0{,}01}=\frac13$.
\end{enumerate}
\resultado convergencia a 0; cota $1/3$.
\comprobacion la desviación típica de $\overline U_{100}$ es $\sqrt{1/300}\approx0{,}058$, de modo que $0{,}1$ son unas $1{,}7$ desviaciones típicas y la cota $1/1{,}7^{2}\approx0{,}33$ es coherente.
\end{solucion}

\begin{solucion}{eje:6-encuesta}
\datos $\hat p$ = media de $n$ variables $\Ber(p)$ independientes.
\planteamiento \cref{cor:t6-frecuencias}.
\pasos
\begin{enumerate}[label=\arabic*.]
  \item (a) $\E[\hat p]=p$ y $\Var(\hat p)=p(1-p)/n$. Por Chebyshev, $\PP(\abs{\hat p-p}>c)\le\PP(\abs{\hat p-p}\ge c)\le\frac{p(1-p)}{nc^{2}}$.
  \item $p(1-p)=\frac14-(p-\frac12)^{2}\le\frac14$, luego la cota es $\le\frac1{4nc^{2}}$.
  \item (b) $\frac1{4n\cdot0{,}05^{2}}\le0{,}05\iff n\ge\frac{1}{4\cdot0{,}0025\cdot0{,}05}=2000$.
\end{enumerate}
\resultado la desigualdad del enunciado; $n=2000$ basta.
\comprobacion (ampliación) la aproximación normal del teorema central del límite, que se estudia en el módulo~5, permite reducir el tamaño a $n\approx(1{,}96/(2\cdot0{,}05))^{2}\approx385$; Chebyshev es más conservador porque vale para cualquier distribución.
\end{solucion}

\begin{solucion}{eje:6-cheb-exacta}
\datos $\mu$, $\sigma>0$, $c\ge\sigma$.
\planteamiento colocar la masa en $\mu-c$, $\mu$ y $\mu+c$, de modo simétrico.
\pasos
\begin{enumerate}[label=\arabic*.]
  \item Sea $\PP(X=\mu\pm c)=\frac{\sigma^{2}}{2c^{2}}$ cada uno y $\PP(X=\mu)=1-\frac{\sigma^{2}}{c^{2}}$. Es una ley válida porque $\sigma^{2}/c^{2}\le1$.
  \item $\E[X]=\mu$ por simetría.
  \item $\Var(X)=c^{2}\cdot\frac{\sigma^{2}}{2c^{2}}\cdot2=\sigma^{2}$.
  \item $\PP(\abs{X-\mu}\ge c)=\PP(X=\mu\pm c)=\frac{\sigma^{2}}{c^{2}}$: igualdad en Chebyshev.
\end{enumerate}
\resultado la variable construida alcanza la cota; ninguna cota que dependa solo de $\sigma$ y $c$ puede ser menor.
\comprobacion con $\sigma=1$, $c=2$: masas $\frac18,\frac34,\frac18$; $\PP(\abs{X-\mu}\ge2)=\frac14=\frac1{2^{2}}$.
\end{solucion}

\begin{solucion}{eje:6-cantelli}
\datos media $\mu$, varianza $\sigma^{2}$, $a>0$.
\planteamiento desplazar por $t\ge0$, elevar al cuadrado, aplicar Markov y optimizar en $t$.
\pasos
\begin{enumerate}[label=\arabic*.]
  \item Si $X-\mu\ge a$, entonces $X-\mu+t\ge a+t>0$ y $(X-\mu+t)^{2}\ge(a+t)^{2}$.
  \item Markov: $\PP(X-\mu\ge a)\le\frac{\E[(X-\mu+t)^{2}]}{(a+t)^{2}}=\frac{\sigma^{2}+t^{2}}{(a+t)^{2}}$, porque $\E[X-\mu]=0$.
  \item Sea $h(t)=\frac{\sigma^{2}+t^{2}}{(a+t)^{2}}$. $h'(t)=\frac{2t(a+t)-2(\sigma^{2}+t^{2})}{(a+t)^{3}}=\frac{2(at-\sigma^{2})}{(a+t)^{3}}$, que se anula en $t^{*}=\sigma^{2}/a$, mínimo de $h$ en $[0,\infty)$.
  \item $h(t^{*})=\frac{\sigma^{2}+\sigma^{4}/a^{2}}{(a+\sigma^{2}/a)^{2}}=\frac{\sigma^{2}(a^{2}+\sigma^{2})/a^{2}}{(a^{2}+\sigma^{2})^{2}/a^{2}}=\frac{\sigma^{2}}{\sigma^{2}+a^{2}}$.
\end{enumerate}
\resultado $\PP(X-\mu\ge a)\le\sigma^{2}/(\sigma^{2}+a^{2})$.
\comprobacion con $a=\sigma$ la cota es $1/2$, mientras que Chebyshev bilateral solo daría $\PP(\abs{X-\mu}\ge\sigma)\le1$.
\end{solucion}

\begin{solucion}{eje:6-expo}
\datos $f(x)=\lambda e^{-\lambda x}$, $x\ge0$.
\planteamiento integración por partes con $u=x^{k}$, $dv=\lambda e^{-\lambda x}dx$, $v=-e^{-\lambda x}$; los términos de borde se anulan porque $x^{k}e^{-\lambda x}\to0$ cuando $x\to\infty$.
\pasos
\begin{enumerate}[label=\arabic*.]
  \item $\E[X]=\bigl[-xe^{-\lambda x}\bigr]_0^\infty+\int_0^\infty e^{-\lambda x}dx=0+\frac1\lambda$.
  \item $\E[X^{2}]=\bigl[-x^{2}e^{-\lambda x}\bigr]_0^\infty+\int_0^\infty2xe^{-\lambda x}dx=\frac2\lambda\E[X]=\frac2{\lambda^{2}}$.
  \item $\Var(X)=\frac2{\lambda^{2}}-\frac1{\lambda^{2}}=\frac1{\lambda^{2}}$.
  \item $\PP(\lambda X>t)=\PP(X>t/\lambda)=e^{-t}$ para $t\ge0$: $\lambda X\sim\Expo(1)$.
\end{enumerate}
\resultado $\E[X]=1/\lambda$, $\Var(X)=1/\lambda^{2}$ y $\lambda X\sim\Expo(1)$.
\comprobacion por el último apartado, $\E[X]=\frac1\lambda\E[\lambda X]=\frac1\lambda\cdot1$, coherente con el primero.
\end{solucion}

\begin{solucion}{eje:6-afin-normal}
\datos $X\sim\Nor(\mu,\sigma^{2})$, $Y=aX+b$, $a\neq0$.
\planteamiento calcular $F_Y$ y reconocer la función de distribución de una normal, $\Phi((y-m)/s)$ (\cref{teo:t6-tipificacion}).
\pasos
\begin{enumerate}[label=\arabic*.]
  \item $a>0$: $F_Y(y)=\PP\bigl(X\le\frac{y-b}a\bigr)=\Phi\Bigl(\frac{(y-b)/a-\mu}{\sigma}\Bigr)=\Phi\Bigl(\frac{y-(a\mu+b)}{a\sigma}\Bigr)$.
  \item $a<0$: al dividir por $a$ se invierte la desigualdad: $F_Y(y)=\PP\bigl(X\ge\frac{y-b}a\bigr)=1-\Phi\Bigl(\frac{(y-b)/a-\mu}\sigma\Bigr)=\Phi\Bigl(\frac{y-(a\mu+b)}{\abs a\sigma}\Bigr)$, usando $1-\Phi(z)=\Phi(-z)$ y $-\frac{(y-b)/a-\mu}{\sigma}=\frac{y-b-a\mu}{-a\sigma}$.
  \item En ambos casos $F_Y(y)=\Phi\bigl((y-m)/s\bigr)$ con $m=a\mu+b$ y $s=\abs a\sigma$: $Y\sim\Nor(a\mu+b,a^{2}\sigma^{2})$.
  \item Con $a=-1$, $b=0$ y $X=Z\sim\Nor(0,1)$: $-Z\sim\Nor(0,1)$.
\end{enumerate}
\resultado $aX+b\sim\Nor(a\mu+b,\,a^{2}\sigma^{2})$.
\comprobacion la media y la varianza coinciden con las que da la linealidad: $\E[Y]=a\mu+b$ y $\Var(Y)=a^{2}\sigma^{2}$.
\end{solucion}
